%% Dokumentenklasse (Koma Script) -----------------------------------------
\documentclass[%
   %draft,     % Entwurfsstadium
   final,      % fertiges Dokument
	 % --- Paper Settings ---
   paper=a4,% [Todo: add alternatives]
   paper=portrait, % landscape
   pagesize=auto, % driver
   % --- Base Font Size ---
   fontsize=12pt,%
	 % --- Koma Script Version ---
   version=last, %
 ]{scrbook} % Classes: scrartcl, scrreprt, scrbook

% Encoding der Dateien (sonst funktionieren Umlaute nicht)
\usepackage{listings}
\usepackage{booktabs}
%\usepackage{acro}
%\usepackage[ruled,vlined]{algorithm2e}
\usepackage[utf8]{inputenc}
\usepackage[htt]{hyphenat}

%%% Preambel
%%% === Textbody ==============================================================
\KOMAoptions{%
   DIV=15,% (Size of Text Body, higher values = greater textbody)
   % DIV=calc % (also areaset/classic/current/default/last) 
   % -> after setting of spacing necessary!   
   BCOR=5mm% (Bindekorrektur)
}%
%%% === Headings ==============================================================
\KOMAoptions{%
   %%%% headings
	% headings=small,  % Small Font Size, thin spacing above and below
   % headings=normal, % Medium Font Size, medium spacing above and below
   headings=big, % Big Font Size, large spacing above and below
   %
   headings=noappendixprefix, % chapter in appendix as in body text
	headings=nochapterprefix,  % no prefix at chapters
   % headings=appendixprefix,   % inverse of 'noappendixprefix'
   % headings=chapterprefix,    % inverse of 'nochapterprefix'
   % headings=openany,   % Chapters start at any side
   % headings=openleft,  % Chapters start at left side
   headings=openright, % Chapters start at right side      
   %%% Add/Dont/Auto Dot behind section numbers 
   %%% (see DUDEN as reference)
   % numbers=autoenddot
   % numbers=enddot
   numbers=noenddot
   % secnumdepth=3 % depth of sections numbering (???)
}%
\setcounter{secnumdepth}{3}
%%% === Page Layout ===========================================================
\KOMAoptions{% (most options are for package typearea)
   twoside=true, % two side layout (alternating margins, standard in books)
   % twoside=false, % single side layout 
   % twoside=semi,  % two side layout (non alternating margins!)
   %
   twocolumn=false, % (true)
   %
   headinclude=false,%
   footinclude=false,%
   mpinclude=true,%      
   %
   headlines=2.1,%
	% headheight=2em,%
   headsepline=true,%
   footsepline=false,%
   cleardoublepage=empty %plain, headings
}%
%%% === Paragraph Separation ==================================================
\KOMAoptions{%
	% parskip=relative, % change indentation according to fontsize (recommanded)
   parskip=absolute, % do not change indentation according to fontsize
   parskip=false    % indentation of 1em
   % parskip=true   % parksip of 1 line - with free space in last line of 1em
   % parskip=full-  % parksip of 1 line - no adjustment
   % parskip=full+  % parksip of 1 line - with free space in last line of 1/4
   % parskip=full*  % parksip of 1 line - with free space in last line of 1/3
   % parskip=half   % parksip of 1/2 line - with free space in last line of 1em
   % parskip=half-  % parksip of 1/2 line - no adjustment
   % parskip=half+  % parksip of 1/2 line - with free space in last line of 1/3
   % parskip=half*  % parksip of 1/2 line - with free space in last line of 1em
}%
%%% === Table of Contents =====================================================
\setcounter{tocdepth}{3} % Depth of TOC Display
\KOMAoptions{%
   %%% Setting of 'Style' and 'Content' of TOC
   % toc=left, %
   toc=indented,%
   %
   toc=bib,
   %toc=nobib,
   % toc=bibnumbered,
   %
	% toc=index,%
   toc=noindex,
   %
   % toc=listof,
   toc=nolistof
   % toc=listofnumbered,
   %   
}%  
%%% === Lists of figures, tables etc. =========================================
\KOMAoptions{%
   %%% Setting of 'Style' and 'Content' of Lists 
   %%% (figures, tables etc)
	% --- General List Style ---
   listof=left, % tabular styles
   listof=indented, % hierarchical style
   % --- chapter highlighting ---
   % listof=chapterentry, % ??? Chapter starts are marked in figure/table
   % listof=chaptergapline, % New chapter starts are marked by a gap 
      		  			   	 % of a single line
	listof=chaptergapsmall, % New chapter starts are marked by a gap 
   	    					   % of a smallsingle line
   % listof=nochaptergap, % No Gap between chapters
   %
   % listof=leveldown, % lists are moved one level down ???
   % --- Appearance of Lists in TOC
    listof=notoc % Lists are not part of the TOC
   %listof=totoc % add Lists to TOC without number
   % listof=totocnumbered, % add Lists to TOC with number
}%  
%%% === Bibliography ==========================================================
%% Setting of 'Style' and 'Content' of Bibliography
\KOMAoptions{%
	% bibliography=oldstyle,%
   bibliography=openstyle,%
   bibliography=nottotoc % Bibliography is not part of the TOC
   % bibliography=totocnumbered, % add Bibliography to TOC with number
   %bibliography=totoc % add Bibliography to TOC without number
}%
%%% === Index =================================================================
%% Setting of 'Style' and 'Content' of Index in TOC
\KOMAoptions{%
   index=nottotoc % index is not part of the TOC
	% index=totoc, % add index to TOC without number
}%
%%% === Titlepage =============================================================
\KOMAoptions{%
   titlepage=true %
   %titlepage=false %
}%
%%% === Miscellaneous =========================================================
\KOMAoptions{% 	
   footnotes=multiple% nomultiple
   %open=any,%
   %open=left,%
   %open=right,%
   %chapterprefix=false,%
   %appendixprefix=false,%
   %chapteratlists=10pt,% entry
}%

\gdef\thetitle{}
\gdef\thethesis{}
\gdef\theauthor{}
\gdef\theabstract{}
\gdef\theabgabedatum{}
\gdef\thegeburtsort{}
\gdef\thegeburtsdatum{}
\gdef\thematrikel{}
\gdef\thesupervisor{}
\gdef\thezusammenfassung{}
%
\def\title#1{\gdef\@title{#1}\gdef\thetitle{#1}}
\def\@title{\@latex@error{No \noexpand\title given}\@ehc}
\def\thesis#1{\gdef\thethesis{#1}}
\def\author#1{\gdef\@author{#1}\gdef\theauthor{#1}}
\def\@author{\@latex@warning@no@line{No \noexpand\author given}}
\def\abgabedatum#1{\gdef\theabgabedatum{#1}}
\def\geburtsort#1{\gdef\thegeburtsort{#1}}
\def\geburtsdatum#1{\gdef\thegeburtsdatum{#1}}
\def\abstract#1{\gdef\theabstract{#1}}
\def\matrikel#1{\gdef\thematrikel{#1}}
\def\supervisor#1{\gdef\thesupervisor{#1}}
\def\zusammenfassung#1{\gdef\thezusammenfassung{#1}}
\input{preambel/preambel}
%
%%%% Neue Befehle
\def\bzgl{bzgl.\ }
\def\dh{d.\,h.\ }
\def\eV{e.\,V.\ }
\def\engl{engl.:\ }
\def\zB{z.\,B.\ }
\def\oae{o.\,\"{a}.\ }
\def\uU{u.\,U.\ }
\def\ua{u.\,a.\ }
\def\ia{i.\,a.\ }
\def\og{o.\,g.\ }
\def\bzw{bzw.\ }
\def\vgl{vgl.\ }
\def\etc{etc.\ }
\def\evtl{evtl.\ }
\def\ca{ca.\ }
\def\sog{sog.\ }
\def\insbes{insbes.\ }
\def\usw{usw.\ }

%--------------------------- References ------------------------------
\def\Eq#1{\eqref{#1}}
\def\Sec#1{Abschnitt~\ref{#1}}
\def\SubSec#1{Subsection~\ref{#1}}
\def\Ex#1{Beispiel~\ref{#1}}
\def\Chap#1{Kapitel~\ref{#1}}
\def\Part#1{Teil~\ref{#1}}
\def\Fig#1{Abbildung~\ref{#1}}
\def\Tab#1{Tabelle~\ref{#1}}
\def\Theo#1{Satz~\ref{#1}}
\def\Def#1{Definition~\ref{#1}}
\def\Alg#1{Algorithmus~\ref{#1}}
\def\Cor#1{Korollar~\ref{#1}}
\def\Lem#1{Lemma~\ref{#1}}
\def\Ax#1{Axiom~\ref{#1}}
\def\Page#1{Seite~\pageref{#1}}

%-------------------------- Debugging ----------------------------------
\def\work#1{#1}
\def\draft#1{}
\def\final#1{}
\def\todo#1{\work{\comment{TODO: #1}}}

%%%%%%%%%%%%%%%%%%%%%%%%%%%%%%%%%%%%%%%%%%%%%%%%%%%%%%%%%%%%%%
%%
%%   Identification
%%
%%%%%%%%%%%%%%%%%%%%%%%%%%%%%%%%%%%%%%%%%%%%%%%%%%%%%%%%%%%%%%
%\NeedsTeXFormat{LaTeX2e}

%%%%%%%%%%%%%%%%%%%%%%%%%%%%%%%%%%%%%%%%%%%%%%%%%%%%%%%%%%%%%%
%%
%%   Preliminary declarations
%%
%%%%%%%%%%%%%%%%%%%%%%%%%%%%%%%%%%%%%%%%%%%%%%%%%%%%%%%%%%%%%%
%\RequirePackage{bm} % correct bold symbols, like \bm
%\RequirePackage{amsmath}
%\RequirePackage{amsfonts}
%\RequirePackage{amssymb}
%\RequirePackage{mathrsfs}
%\RequirePackage{dsfont} % double stroke symbols like IR
%\RequirePackage[single]{accents} % own accents 
%\RequirePackage{exscale} % Scaling of Math Symbols?

%%%%%%%%%%%%%%%%%%%%%%%%%%%%%%%%%%%%%%%%%%%%%%%%%%%%%%%%%%%%%%
%%
%%   Options
%%
%%%%%%%%%%%%%%%%%%%%%%%%%%%%%%%%%%%%%%%%%%%%%%%%%%%%%%%%%%%%%%



%%%%%%%%%%%%%%%%%%%%%%%%%%%%%%%%%%%%%%%%%%%%%%%%%%%%%%%%%%%%%%
%%
%%   More declarations
%%
%%%%%%%%%%%%%%%%%%%%%%%%%%%%%%%%%%%%%%%%%%%%%%%%%%%%%%%%%%%%%%


%------------------------------------------------------------
% Mathematical Formulae
%------------------------------------------------------------
\newcommand{\rv}[1]{\ensuremath{{\boldsymbol{#1}}}}
\renewcommand{\vec}[1]{\ensuremath{{\underline{#1}}}}
\newcommand{\rvec}[1]{\ensuremath{{\boldsymbol{\underline{#1}}}}}
\newcommand{\mat}[1]{{\ensuremath{{\mathbf{#1}}}}}

\newcommand{\Menge}[1]{{\mathcal{#1}}}
\newcommand{\SquareRoot}[1]{{\mbox{\bfseries\sffamily #1}}}
\newcommand{\MatrixMenge}[2]{{\mathcal{M}}_{#1#2}}
\newcommand{\intndim}{\ds\int_{-\infty}^\infty \dots \int_{-\infty}^\infty}
\newcommand{\Eye}{\mat{I}}

% operators (in alphabetical order):
\renewcommand{\arg}{\operatorname*{arg}}
%\DeclareMathOperator*{\argmin}{arg\,min}
%\DeclareMathOperator*{\argmax}{arg\,max}
\newcommand{\argmin}[1]{\mathop{\mathrm{arg}\; \min_{#1}}}
\newcommand{\argmax}[1]{\mathop{\mathrm{arg}\; \max_{#1}}}
\DeclareMathOperator{\col}{col}
\newcommand{\dd}{\operatorname{d}\!}
%\def\d{\ensuremath{\operatorname{d}}} % for \int f(x)\d x or \d y / \d y
%\renewcommand{\d}{\ensuremath{\text{d}}}
\DeclareMathOperator{\diag}{diag}
\DeclareMathOperator{\Cov}{Cov}
\DeclareMathOperator{\erf}{erf}
\DeclareMathOperator{\E}{E}
\DeclareMathOperator{\Proj}{proj}
\renewcommand{\Pr}{\operatorname{Pr}} 
\DeclareMathOperator{\rang}{rang}
\DeclareMathOperator{\rank}{rank}
\DeclareMathOperator{\jj}{j}% complex unit
\DeclareMathOperator{\rect}{rect}
\def\T{^\mathrm{T}} % transpose
\DeclareMathOperator{\trace}{trace}
\DeclareMathOperator*{\Var}{Var} 
\renewcommand{\arctan}{\operatorname{arctan}}
\newcommand{\atann}{\operatorname{arctan2}}
%\DeclareMathOperator{\atan}{atan}
\def\atan2{\operatorname{atan2}}
\DeclareMathOperator{\pot}{pot}

\newcommand{\MatrixTyp}[2]{\Menge{M}_{#1, #2}}
\newcommand{\FromTo}[2]{$#1 = 1,\, \ldots,\, #2$}
\newcommand{\LMInd}{{\operatorname{LM}}}

%---------------------------------------------------------------------
% Abbreviations
%---------------------------------------------------------------------

\def\({\left(}
\def\){\right)}
\def\[{\left[}
\def\]{\right]}

\newcommand{\ds}{\displaystyle}
\newcommand{\ns}{\textstyle}

\newcommand{\NewB}{\ensuremath{\mathds{R}}}
\newcommand{\NewC}{\ensuremath{\mathds{C}}}
\newcommand{\NewN}{\ensuremath{\mathds{N}}}
\newcommand{\NewQ}{\ensuremath{\mathds{Q}}}
\newcommand{\NewR}{\ensuremath{\mathds{R}}}
\newcommand{\NewZ}{\ensuremath{\mathds{Z}}}
\newcommand{\IR}{\NewR}
\newcommand{\IN}{\NewN}
\newcommand{\IZ}{\NewZ}
\newcommand{\IC}{\NewC}

\newcommand{\Gauss}{{\mathcal{N}}}

\setcounter{MaxMatrixCols}{20} %Fuer ams umgebung


%--------------------------------------------------------------------------
%Einheiten mit korrekten Leerzeichen
%Usage: \Meter{\micro}{10} results $10\mu m$ 
\newcommand{\Meter}[2][]{\unit{#2}#1\meter} 
\newcommand{\Volt}[2][]{\unit{#2}#1\volt} 
\newcommand{\Hertz}[2][]{\unit{#2}#1\hertz} 
\newcommand{\Second}[2][]{\unit{#2}#1\second} 
\newcommand{\Pascal}[2][]{\unit{#2}#1\pascal} 
\newcommand{\Euro}[2][]{\unit{#2}#1\euro}

\newcommand{\nm}[1] {\Meter[\nano]{#1}}
\newcommand{\mum}[1]{\Meter[\micro]{#1}}
\newcommand{\mm}[1] {\Meter[\milli]{#1}}
\newcommand{\cm}[1] {\Meter[\centi]{#1}}
\newcommand{\mV}[1] {\Volt[\milli]{#1}}
\newcommand{\ms}[1] {\Second[\milli]{#1}}
\newcommand{\KHz}[1] {\Hertz[\kilo]{#1}}

%--------------------------- Operators ---------------------------
\def\TJ{(TJ)}
\def\TkJ{(T^kJ)}
\def\TJp{(T_\pi J)}
\def\TkJp{(T_\pi^k J)}
\def\vTJ{T  \vec J}
\def\vTkJ{T^k\vec J}
\def\vTJp{T_\pi \vec J}
\def\vTkJp{T_\pi^k \vec J}

%------------------------ Random variables ---------------------------
\def\v{\rv v}
\def\w{\rv w}
\def\x{\rv x}
\def\y{\rv y}
\def\z{\rv z}

%---------------------------- Vectors --------------------------------
\def\vdel{\vec \delta}
\def\veta{\vec \eta}
\def\vOmega{\vec \Omega}
\def\vsig{\vec \sigma}
\def\va{\vec a}
\def\vb{\vec b}
\def\vc{\vec c}
\def\ve{\vec e}
\def\vg{\vec g}
\def\vh{\vec h}
\def\vmu{\vec \mu}
\def\vn{\vec n}
\def\vpi{\vec \pi}
\def\vu{\vec u}
\def\vv{\vec v}
\def\vw{\vec w}
\def\vx{\vec x}
\def\vy{\vec y}
\def\vz{\vec z}
\def\vxi{\vec \xi}
\def\vB{\vec B}
\def\vF{\vec F}
\def\vJ{\vec J}
\def\vM{\vec M}
\def\vP{\vec P}
\def\vX{\vec X}

%------------------------- Random vectors ----------------------------
\def\rvn{\rvec n}
\def\rvv{\rvec v}
\def\rvw{\rvec w}
\def\rvx{\rvec x}
\def\rvy{\rvec y}
\def\rvz{\rvec z}
\def\rvhz{\rvec{\hat z}}
\def\rvxi{\rvec \xi}
\def\rvX{\rvec X}

%---------------------------- Matrices -------------------------------
\def\A{\mat A}
\def\B{\mat B}
\def\C{\mat C}
\def\D{\mat D}
\def\mE{\mat E}
\def\F{\mat F}
\def\H{\mat H}
\def\I{\mat I}
\def\K{\mat K}
\def\L{\mat L}
\def\M{\mat M}
\def\P{\mat P}
\def\Q{\mat Q}
\def\R{\mat R}
\def\V{\mat V}
\def\W{\mat W}

%------------------------------ Mean ---------------------------------
\def\ha{\hat a}
\def\hv{\hat v}
\def\hw{\hat w}
\def\hx{\hat x}
\def\hy{\hat y}
\def\hz{\hat z}
\def\hva{\hat{\vec a}}
\def\hvv{\hat{\vec v}}
\def\hvw{\hat{\vec w}}
\def\hvx{\hat{\vec x}}
\def\hvy{\hat{\vec y}}
\def\hvz{\hat{\vec z}}
\def\hvX{\hat{\vec X}}

%------------------------------ Sets ---------------------------------
\def\SA{\mathcal A}
\def\SB{\mathcal B}
\def\SE{\mathcal E}
\def\SF{\mathcal F}
\def\SG{\mathcal G}
\def\SI{\mathcal I}
\def\SL{\mathcal L}
\def\SM{\mathcal M}
\def\SN{\mathcal N}
\def\SO{\mathcal O}
\def\SP{\mathcal P}
\def\SS{\mathcal S}
\def\SU{\mathcal U}
\def\SV{\mathcal V}
\def\SW{\mathcal W}
\def\SX{\mathcal X}
\def\SY{\mathcal Y}
\def\SZ{\mathcal Z}

%------------------------------ Misc ---------------------------------
\def\ii{^{(i)}}
\def\ij{^{(j)}}
\def\G{G(\veta, \gamma)}

%--------------------------- Densities -------------------------------
\def\fc{\tilde f(y|x)}
\def\fvc{\tilde f(\vy|\vx)}
\def\fcxy{f(x|y)}
\def\fcvxy{f(\vx|\vy)}
\def\fcyx{f(y|x)}
\def\fcvyx{f(\vy|\vx)}
\def\fija{f_{ij}(a)}
\def\fijak{f_{ij}^k(a)}
\def\fh{f(y, x; \veta)}
\def\fvh{f(\vy, \vx; \veta)}
\def\fkex{f_k^e(x_k)}
\def\fkevx{f_k^e(\vx_k)}
\def\fkpx{f_{k}^p(x_{k})}
\def\fkkpx{f_{k+1}^p(x_{k+1})}
\def\fkpvx{f_{k}^p(\vx_{k})}
\def\fkkpvx{f_{k+1}^p(\vx_{k+1})}
\def\ftxa{f(x_{k+1}|x_k, a_k)}
\def\ftvxa{f(\vx_{k+1}|\vx_k, \va_k)}
\def\fkv{f_k^v(v_k)}
\def\fkvv{f_k^v(\vv_k)}
\def\fkw{f_k^w(w_k)}
\def\fkvw{f_k^w(\vw_k)}
\def\fnex{f_n^e(x_n)}
\def\fnevx{f_n^e(\vx_n)}
\def\fnpx{f_{n}^p(x_{n})}
\def\fnnpx{f_{n+1}^p(x_{n+1})}
\def\fnpvx{f_{n}^p(\vx_{n})}
\def\fnnpx{f_{n+1}^p(\vx_{n+1})}
\def\fx{f(x)}
\def\fxy{f(x,y)}
\def\fvx{f(\vx)}
\def\fvxy{f(\vx,\vy)}
\def\fy{f(y)}
\def\fvy{f(\vy)}

%--------------------------- Distributions ---------------------------
\def\Fc{\tilde{F}(y|x)}
\def\Fh{F(y, x; \veta)}

%------------------------------ Theorems ------------------------------------
\theoremheaderfont{\normalfont\bfseries\sffamily}
\theoremstyle{plain}
{\theorembodyfont{\itshape} \newtheorem{theorem}     {Satz}[chapter]}
{\theorembodyfont{\itshape} \newtheorem{lemma}       {Lemma}[chapter]}
{\theorembodyfont{\itshape}\newtheorem{corollary}   {Korollar}[chapter]}
%{\theorembodyfont{\sffamily}\newtheorem{example}     {Beispiel}[chapter]}
{\theorembodyfont{\normalfont} \newtheorem{remark}      {\sffamily Remark}[chapter]}
{\theorembodyfont{\itshape}\newtheorem{definition}  {Definition}[chapter]}
{\theorembodyfont{\itshape}\newtheorem{axiom}  {Axiom}[chapter]}
%\theoremheaderfont{\scshape\normalsize}
%\theoremstyle{plain}
%{\theorembodyfont{\normalfont\small} \newtheorem{Proof}	{Proof}}
\newenvironment{proof}
{\noindent{\scshape Beweis.}}
{\hspace*{\fill}$\square$\bigskip}

\makeatletter 
\def\ifEmpty#1{\def\@temp{#1}\ifx\@temp\@empty} 
\makeatother

\newcounter{example}[chapter]
\renewcommand{\theexample}{\thechapter.\arabic{example}}
\newenvironment{example}[1][]
{\refstepcounter{example}\par\bigskip\noindent%
   {\normalfont\bfseries\sffamily Beispiel~\theexample~\ifEmpty{#1}\else(#1)\fi} \sffamily}
{\hspace*{\fill}$\blacksquare$\smallskip\smallskip}

\newenvironment{principle}[1][]
{\par\bigskip\noindent%
   {\normalfont\bfseries\sffamily \ifEmpty{#1}\else#1\fi} \itshape}
{\smallskip\smallskip}
\algrenewcommand{\algorithmiccomment}[1]{\hfill// #1}

\newcommand{\emphc}[1]{\comment{\normalfont\marginfont #1}\emph{#1}}

% Hervorheben & Index
\newcommand{\emphi}[2][]{%
	\emph{#2}%
	\ifthenelse{\isempty{#1}}{\index{#2}}{\index{#1}}%
}

% Hervorheben, Index & Marginalkommentar
\newcommand{\emphci}[2][]{%
	\comment{\normalfont\marginfont #2}%
	\emph{#2}%
	\ifthenelse{\isempty{#1}}{\index{#2}}{\index{#1}}%
}

\usetikzlibrary{topaths}
\usetikzlibrary{decorations.pathmorphing}
\usetikzlibrary{fit}

\newcommand{\RA}{\tikz[baseline=0mm]  \draw[ultra thick, -stealth] (0,1.3mm)--(3.5mm,1.3mm); }

\newcommand{\coordcross}[6]{%
	\draw[->, >=latex] (#1,0) -- (#2,0) node[right,below] {#5};
  \draw[->, >=latex] (0,#3) -- (0,#4) node[left] {#6};
}

%\newcommand{\xtick}[3]{%
%}

\tikzstyle{blockdef}=[rectangle, rounded corners, draw=black, thick, text centered]
\tikzstyle{rblock}=[blockdef, inner sep=6pt, text width=2.7cm]
\tikzstyle{cblock}=[blockdef, circle, inner sep=2pt, text width=2mm]
\tikzstyle{ablock}=[ultra thick, -latex, black]
\tikzstyle{aablock}=[ultra thick, latex-latex, black]
\tikzstyle{fblock}=[draw=black, fill]
\tikzstyle{fitblock}=[blockdef, fill=none, dashed, inner sep = 1mm]


%% Color Boxes
\tikzstyle{hlNone}=[anchor=base, rounded corners=3pt]
\tikzstyle{hlFrameNone}=[anchor=base, rounded corners=3pt, thick]
\tikzstyle{hlFrameBlue}=[hlFrameNone, fill=BlockBlue, draw=HilightBlue]
\tikzstyle{hlFrameRed}=[hlFrameNone, fill=BlockRed, draw=HilightRed]
\tikzstyle{hlBlue}=[hlNone, fill=BlockBlue]
\tikzstyle{hlGreen}=[hlNone, fill=BlockGreen]
\tikzstyle{hlLightGreen}=[hlNone, fill=BlockGreen!40]
\tikzstyle{hlRed}=[hlNone, fill=BlockRed]

\everymath{\displaystyle}
\def\hlFrameBlue<#1>#2{\alt<#1>{
  \tikz[baseline]{\node[hlFrameBlue] {#2};}
}{
  \tikz[baseline]{\node[hlFrameNone] {#2};}
}}

\def\hlFrameRed<#1>#2{\alt<#1>{
  \tikz[baseline]{\node[hlFrameRed] {#2};}
}{
  \tikz[baseline]{\node[hlFrameNone] {#2};}
}}

\def\hlBoxBlue<#1>#2{\alt<#1>{
  \tikz[baseline]{\node[hlBlue] {#2};}
}{
  \tikz[baseline]{\node[hlFrameNone] {#2};}
}}


\def\hlBoxGreen<#1>#2{\alt<#1>{
  \tikz[baseline]{\node[hlGreen] {#2};}
}{
  \tikz[baseline]{\node[hlFrameNone] {#2};}
}}

\def\hlBoxLightGreen<#1>#2{\alt<#1>{
  \tikz[baseline]{\node[hlLightGreen] {#2};}
}{
  \tikz[baseline]{\node[hlFrameNone] {#2};}
}}

\def\hlBoxRed<#1>#2{\alt<#1>{
  \tikz[baseline]{\node[hlRed] {#2};}
}{
  \tikz[baseline]{\node[hlFrameNone] {#2};}
}}
%%% Silbentrennung
\hyphenation{Strategie-iteration}
\hyphenation{Werte-iteration}
\hyphenation{Sicherheits-�quivalenz}

%%% Wasserzeichen "Entwurf"
\draft{%
\AddToShipoutPicture{%
  \setlength{\unitlength}{1cm}
  %\ifthenelse{\equal{\value{page}}{1}}{}{%
  \ifthenelse{\isodd{\value{page}}}
  	{\put(19,14.98){% Mitte rechts von A4 (21cm max 29.7cm)
      \makebox(0,0){\rotatebox{90}{\textcolor[gray]{0.8}{\fontsize{4cm}{4cm}\selectfont{\rm{Entwurf}}}}}
    }}
    {\put(2,14.98){ % Mitte links von A4 (21cm max 29.7cm)
        	\makebox(0,0){\rotatebox{-90}{\textcolor[gray]{0.8}{\fontsize{4cm}{4cm}\selectfont{\rm{Entwurf}}}}}
    }}
  }
%}
}

%%% Design für listings
\definecolor{codegreen}{rgb}{0,0.6,0}
\definecolor{codegray}{rgb}{0.5,0.5,0.5}
\definecolor{codepurple}{rgb}{0.58,0,0.82}
\definecolor{backcolour}{rgb}{0.95,0.95,0.92}

\lstdefinestyle{codestyle}{
    backgroundcolor=\color{backcolour},   
    commentstyle=\color{codegreen},
    keywordstyle=\color{magenta},
    numberstyle=\tiny\color{codegray},
    stringstyle=\color{codepurple},
    basicstyle=\ttfamily\footnotesize,
    breakatwhitespace=false,         
    breaklines=true,                 
    captionpos=b,                    
    keepspaces=true,                 
    numbers=left,                    
    numbersep=5pt,                  
    showspaces=false,                
    showstringspaces=false,
    showtabs=false,                  
    tabsize=2,
    frame=lines
}

\lstset{style=codestyle}


%%% Formatierung
%\clubpenalty=1000000
%\widowpenalty=1000000
%\brokenpenalty=1000000
%\finalhyphendemerits=1000000
%\renewcommand{\baselinestretch}{1.7}

%% Dokument Beginn %%%%%%%%%%%%%%%%%%%%%%%%%%%%%%%%%%%%%%%%%%%%%%%%%%%%%%%%

% - Deckblatt,
% - Inhaltsverzeichnis,
% - Hauptteil gegliedert z.B. in
%   Einleitung, Grundlagen, Experimente, Ergebnisse, Zusammenfassung
% - Literaturverzeichnis,
% - Abbildungsverzeichnis (ggf.),
% - Tabellenverzeichnis (ggf.),
% - Abkürzungsverzeichnis (ggf.),
% - Formelverzeichnis (ggf.),
% - Anhang, (nicht mehr Bestandteil der Arbeit! Wird daher nicht bewertet)
% - Erklärung der Urheberschaft,

\title{Automated Deployment of Machine Learning Applications on a Scalable Heterogeneous GPU-Accelerated Apache Spark Cluster}
\author{Marcel Pascal Stolin}
\matrikel{32168}
\thesis{Bachelor Thesis}
\geburtsdatum{03.04.1993}
\geburtsort{Kamen}
\abgabedatum{\today}
\supervisor{M.Sc. Christoph Hennebold}

\zusammenfassung{Die wachsende Rechenkomplexität beim trainieren von Machine Learning-Modellen mit Big Data erfordert eine Erhöhung der Skalierbarkeit von Rechenclustern. In verteilten Systemen kann ein Rechencluster dynamisch skaliert werden, sobald eine Leistungsschwelle erreicht wird. Ein beliebtes Framework für verteiltes Rechnen ist Apache Spark, das die parallele Datenverarbeitung auf einem Rechencluster ermöglicht. Mit der jüngsten Entwicklung der RAPIDS-Plugin-Suite kann die Rechenleistung weiter gesteigert werden, indem Apache Spark die Nutzung von GPUs ermöglicht wird.
Außerdem ist durch die Popularität der DevOps-Kultur die Automatisierung von Entwicklungsprozessen zu einem trivialen Prozess geworden. Das Training von Machine Learning-Modellen ist ein komplexer Prozess, der mithilfe einer CI-Pipeline automatisiert werden kann.
Ziel dieser Arbeit ist es, ein Entwurf und die Implementierung einer skalierbaren verteilten Rechenumgebung für das Training von Machine Learning-Modellen mit Apache Spark vorzustellen. Dazu wird die Leistung des Apache Spark-Clusters durch die Nutzung von GPUs unter Verwendung der RAPIDS-Plugin-Suite erhöht. Ein \textit{Auto-Scaler}, der die Anzahl der Apache Spark-Worker skaliert, wird implementiert und der Rechenumgebung hinzugefügt, um die Skalierbarkeit weiter zu erhöhen. Darüber hinaus wird ein CI-Pipeline-Entwurf und eine Implementierung mittels GitLab CI/CD vorgestellt, die das Training einer Machine Learning-Anwendung auf Apache Spark automatisiert. Die Leistung der Rechenumgebung wird mit zwei unterschiedlichen Machine Learning-Algorithmen bewertet. Die Ergebnisse zeigen, dass die Leistung des Apache Spark-Clusters durch die Verwendung von GPUs signifikant verbessert wurde, während die \textit{Auto-Scaler}-Implementierung die Leistung des Apache Spark-Clusters nicht erhöht hat.\newpage}
\abstract{The growing computational complexity of training Machine Learning models on Big Data requires increasing the scalability of computing clusters. With distributed computing, a cluster can dynamically be scaled whenever a performance threshold is reached. A popular framework for distributed computing is Apache Spark, which enables parallel data processing on a large computing cluster. With the recent development of the RAPIDS plugin suite, the computational power can further be increased by enabling Apache Spark to leverage GPUs. Furthermore, due to DevOps culture popularity, automating development processes has become a trivial process. The training of machine learning models is a complex process that can be automated using a CI pipeline. This thesis aims to present a design and implementation of a scalable distributed computing environment for training Machine Learning models using Apache Spark. Therefore, the Apache Spark cluster's performance is increased by leveraging GPUs using the RAPIDS plugin suite. An \textit{Auto-Scaler} that scales the number of Apache Spark workers is implemented and added to the computing environment to increase the scalability further. Additionally, a CI pipeline design and implementation using GitLab CI/CD is introduced, which automates the deployment of an application to the Apache Spark cluster. The performance of the computing environment is evaluated using two different Machine Learning algorithms. The results show that the performance of the Apache Spark cluster has significantly been improved using GPUs, whereas the \textit{Auto-Scaler} implementation did not increase the performance of the Apache Spark cluster.}

\begin{document}
% Deckblatt
\begin{titlepage}
\begin{minipage}[h]{1.0\linewidth+1.6cm}
  \begin{center}
		\rule[2ex]{\textwidth}{1pt}
		\begin{tabular}{ccc}
		 \includegraphics[width=.4\textwidth]{logo_iff} & \hphantom{abcdefghijklmnop} & \includegraphics[width=.3\textwidth]{logo_ipa}
		\end{tabular}
		\\[3mm]
    \rule[10ex]{\textwidth}{1pt}
		\huge\textbf{\textsf{\thethesis}}\\[2cm]
  	\Large \textbf{\textsc{\thetitle}} \\[1cm]
    \normalsize
    vorgelegt von\\[1ex]
		{\large \textbf{\theauthor}}\\[1ex]
		\thematrikel \\[1ex]
		geboren am \thegeburtsdatum~in~\thegeburtsort \\
    \vspace{2.5cm}
    Verfasst am\\[1ex]
		\textbf{\textsf{Fraunhofer-Institut f�r Produktionstechnik und Automatisierung IPA}}\\[1ex]
		und\\[1ex]
		\textbf{\textsf{Institut f�r Industrielle Fertigung und Fabrikbetrieb,\\Universit�t Stuttgart}}\\
    \vspace{2cm}
  \end{center}
	\normalsize{
    \begin{tabular}{ll}
    	Pr�fer: & Univ.-Prof. Dr.-Ing. habil. Marco Huber \\
      Betreuung: & \thesupervisor \\
    \end{tabular}\\[2ex]
    }
	\rule[-2ex]{\textwidth}{1pt}
  \begin{flushright}
		\footnotesize Stuttgart, \theabgabedatum
	\end{flushright}
\end{minipage}
\end{titlepage}


\frontmatter

\cleardoublepage
\chapter*{Eidesstattliche Erkl�rung}

Ich versichere hiermit, dass ich, \theauthor, die vorliegende \thethesis\ selbstst�ndig angefertigt, keine anderen als die angegebenen Hilfsmittel benutzt und sowohl w�rtliche als auch sinngem�� entlehnte Stellen als solche kenntlich gemacht habe. Die Arbeit hat in gleicher oder �hnlicher Form noch keiner anderen Pr�fungsbeh�rde vorgelegen.

\bigskip\bigskip\bigskip
\noindent Ort, Datum: \underline{\hphantom{01.01.2020, L�denscheidt-West}}

\bigskip\bigskip\bigskip
\noindent Unterschrift: $\underset{\text{\small\theauthor}}{\underline{\hphantom{\text{01.01.2020, L�denscheidt-West}}}}$

\cleardoublepage
\chapter*{Zusammenfassung}
\thezusammenfassung
%
\bigskip\bigskip\bigskip\bigskip\bigskip
\begin{flushright}
	{\LARGE\sffamily \textbf{Abstract}}
\end{flushright}
\bigskip\bigskip\bigskip
\theabstract

\cleardoublepage
% Abbildungs- und Tabellenverzeichnis
\listoffigures
\listoftables
\lstlistoflistings
% Inhaltsverzeichnis in den PDF-Links eintragen
% \pdfbookmark[1]{Inhaltsverzeichnis}{toc}
\tableofcontents

\chapter*{Notation}

\section*{Konventionen}
%
\begin{tabular}{rl}
$x$ & Skalar
\\
$\vx$ & Spaltenvektor
\\
$\x, \rvx$ & Zufallsvariable/-vektor
\\
$\hx, \hvx$ & Mittelwert/-vektor
\\
$x^*, \vx^*$ & Optimaler Wert/Vektor
\\
$x_{0:k}, \vx_{0:k}$ & Folge von Werten $(x_0, x_1, \ldots, x_k)$ / Vektoren $(\vx_0, \vx_1, \ldots, \vx_k)$
\\
$\A$ & Matrizen
\\
$\SA$ & Mengen
\\
$\preceq, \prec$ & schwache/strenge Pr�ferenzrelation
\\
$\NewR$ & Reelle Zahlen
\\
$\NewN$ & Nat�rliche Zahlen
\\
$\blacksquare$ & Ende eines Beispiels
\\
$\square$	& Ende eines Beweises
\end{tabular}

\section*{Operatoren}
%
\begin{tabular}{rl}
$\mat A\T$ & Matrixtransposition
\\
$\mat A^{-1}$ & Matrixinversion
\\
$|\mat A|$ & Determinante einer Matrix
\\
$|\SA|$ & Kardinalit�t der Menge $\SA$
\\
$\pot(\SA)$ & Potenzmenge von $\SA$
\\
$\E\{\cdot\}$ & Erwartungswertoperator
\\
$\SO(g)$ & O-Kalk�l entsprechend der Landau-Notation bei welcher \\&beispielsweise $f(x) \in \SO(g(x))$ besagt, dass die Funktion \\&$f(x)$ die Komplexit�t $\SO(g(x))$ besitzt
\end{tabular}

\section*{Spezielle Funktionen}
%
\begin{longtable}{rl}
$\Pr(\SE)$ & Wahrscheinlichkeitsma�, welches die Wahrscheinlichkeit
\\ & angibt, dass Ereignis $\SE$ eintritt
\\
$p(\vx)$ & (Wahrscheinlichkeits-)Dichtefunktion f�r kontinuierliche $\vx$ \\
& und Z�hldichte f�r diskrete $\vx$
\\
$p(\vx|\vy)$ & Bedingte Dichtefunktion
\\
$P(\vx)$ & (Wahrscheinlichkeits-)Verteilungsfunktion
\\
$\erf(x)$ & Gau�'sche Fehlerfunktion
\\
$\exp(x)$ & Exponentialfunktion $e^x$
\\
$\Gauss(\vx; \hvx, \C_x)$ & Gau�dichte, \dh Dichtefunktion eines normalverteilten \\&Zufallsvektors $\rvx$ mit Mittelwertvektor $\hvx$ und\\& Kovarianzmatrix $\C_x$
\end{longtable}

% Hauptteil
\mainmatter

%\iffalse
\chapter{Introduction}
\label{sec:demo}
%


%Complex computation operations rely on a IT infrastructure that is able to perform operations on scale. Computing systems that automatically adapt to demands and conditions of the workload serve the needs of
%highly scaled computing systems QUELLE. The microservice architecture is a technique to divide a big application into independent services. Each single service only serves one business task of the application QUELLE. In recent years, container applications have become popular to deploy microservices and visualize the needed resources for performing complex operations QUELLE. A container is a single unit that serves the 
%whole environment for a microservice.

% -- Struktur von Intro
% Grid computing
% Autonomic computing
% Microservices ??
% Container
% GPU acceleration


% Start
Storing huge amounts of data has become inexpensive in recent years, but processing it, requires 
parallel computations on clusters with multiple machines \cite{Chambers2018Spark}. Complex 
computation operations rely on a IT infrastructure with the ability to perform operations on scale. 
%Scalability is the ability of a system to grow in accordance of the amount of work \cite{TakadaDistr}.

% Sachen aus autonomic computing
In order to achieve high scalability, computing systems need to adapt dynamically to demands and conditions of the workload.


% Technologies
%In recent years, technologies have contributed to improve the scalability and performance of modern 
%computing systems like cloud computing, virtualization and GPU acceleration.

% Autonomic computing
%Autonomic computing is an architecture for computing systems to manage themselves in accordance to high level objectives, %configured by administrators. These computing systems dynamically adapt to demands and conditions of the workload. The f%our main attributes of autonomic computing are self-configuring, self-healing, self-optimizing and self-protecting \%cite{Kephart2003VisionComputing}. The autonomic computing architecture consists 
%of an autonomic manager, which is responsible to make adjustments in it`s own environment \cite{Sinreich2006AnAB}. %Technologies like Docker help to create containers ... bla bla bla

% Virtualization
%Modern virtualization technologies like containers help to overcome scalability limitations with the creation of virtual instances %of IT resources \cite{Erl2013CloudComputing}.


% GPU acceleration
% In addition to elasticity
%Another aspect to improve complex computation operations is GPU acceleration. GPU accelerations has been an on-going %research area in the recent years.

\section{Problem statement}

%% Struktur
% Möglichkeit IT resourcen automatisch nach bedarf zu skalieren
% Viele ETL operation sind daten intensiv von natur, GPU kann helfen

%To prevent a breakdown of the IT environment while complex operations are being computed, the IT environment should have the ability to scale IT resources dynamically
%when performance peaks and fluctuations occur.

% Dynamisch auf performance spitzen reagieren
ETL \footnote{Extract transform load} operations are compute-intensive. During the execution of analytic applications, performance thresholds can be reached and
the computing system can become out-of-order. 

In addition, many algorithms profit from data-parallelism. % Das ha mehr mit SPark zu tun

%If performance thresholds are reached and the IT environment can`t adapt 
%dynamically to the usage demand, the environment can look temporarily out-of-order. 
%During the execution of operations, performance threshold can be reached and leads to 

% Mit GPU ETLs beschleunigen
%In addition, algorithms do profit from data-parallelism. 

\section{Research questions}

\section{Thesis structure}
\chapter{Theoretical Foundation}
\label{sec:related}
%

This chapter provides background information about ...
\todo{Describe Chapter}

\section{Scalability}
% What is scalability
Scalability defines the ability of a computing system to handle an increasing amount of load \cite{Farcic2017Toolkit21}.
% Why scaling
The main reasing for scaling is High-Availability TOOLKIT 2.1.
% Choose right approach
To choose the right approach of scaling an environment is essential.
% Different kind of scaling
To scale the computing capacity of a system, different approaches exists. A major scaling approach is horizontal scaling.

%% Mehr darüber, warum scaling eigentlich so wichtig ist, und welche probleme man damit löst


\subsection{Horizontal Scaling}
% What is horizontal scaling
Horizontal scaling is accomplished by duplicating nodes in the computing environment \cite{Wilder2012CloudPatterns}.
% DIstribution
Dplucation the nodes in a computing environment increases the computational capacity of the environment. In addition, the workload can be distributed across all clones to handle and balance and increasing load of tasks \cite{Wilder2012CloudPatterns, Abbott2015ScalabilityArt}. 
% What about homogeneous nodes
To increase the efficiency of horizontal scaling, all nodes should be homogeneous. Homogeneous nodes are able to perform the same work and response as other nodes \cite{Abbott2015ScalabilityArt}.


\subsection{Limitations of Scalability}
The limit of scalability is reached, when a computing system is not able to serve the requests of it`s concurrent users \cite{Wilder2012CloudPatterns}.
% hardware capacity
If the scalability of a computing environment is reached, an option is to add more powerful hardware resources to the system. This approach is called vertical scaling QULLE??. By adding more powerful hardware resources, the point can be reached, where more powerful hardware becomes unaffordable or not available \cite{Wilder2012CloudPatterns}.
% Why horizontal i recommended
To overcome the limits of hardware cpacity, a computing system shuld be designed to scale horizontally in the first place \cite{Abbott2015ScalabilityArt}.


% ===========================================
% ===========================================
\section{Self-Healing Environments}
% Auch über container


% ===========================================
% ===========================================
\section{Monitoring}
% Monitoring should be able to deal with dynanism
Self-healing and self-adopting systems are very dynamic in nature. A monitoring should be able to deal with a dynamic environment.
% Requirements of a monitoring system
The requirements for a monitoring system, that is able to monitor a dynamic changing environment, are the following:
\begin{itemize}
\item An efficient database to store metrics
\item A decentralized way of generating metrics \cite{Farcic2017Toolkit21} % Neu schreiben
\item Multi-dimensional metrics \cite{Farcic2017Toolkit21}
\item A powerful query language \cite{Farcic2017Toolkit21}
\end{itemize}


\subsection{Database}
% TS are needed
Saving continous data needs to be done efficiently. 
For time-based metrics, a time series database is he most effective way to save the data \cite{Farcic2017Toolkit21}. To store a huge amount of data, data is stored in a very compact format.


\subsection{Multi-dimensional Metrics}
% Why do we need dimensions
For query languages to be effective, metrics need to be dimensional. Metric without dimensions, are limited in their capabilities.
% Dynamic environment
In a dynamic environment, services are dynamically added and removed. Therefore, a dynamic environment needs dynamic analytics where metrics represent all dimension in the environment \cite{Farcic2018Toolkit22}.
% Example
\begin{lstlisting}[frame=single, label=lst:mon_metr_dimless, caption=Example of a dimensionless-metric, captionpos=b]
container_memory_usage
\end{lstlisting}
\begin{lstlisting}[frame=single, label=lst:mon_metr_withdim, caption=Example of a metric with dimensions, captionpos=b]
container_memory_usage{service="my_service"}
\end{lstlisting}
% Explain
As the exasmple \Lst{lst:mon_metr_dimless} and \Lst{lst:mon_metr_withdim} show, the metric with dimension provide more efficient querieing to gather informations about the environment.


\subsection{Query Model}


\subsection{Choosing a Monitoring Tool}
Graphite or Prometheus, Prometheus it its!

\chapter{Related Work}
\label{chap:03_related-work}


This chapter provides an overview of related literature for this thesis. Furthermore, the surveyed literature is built on the theoretical foundation introduced in \Chap{chap:02_foundation}. This chapter introduces work about auto-scaling computing environments, GPU accelerated Apache Spark cluster, and the implementation of an automated deployment pipeline. These topics are related to the choice of technologies (\Chap{chap:04_background}), the proposed conceptual design of this thesis (\Chap{chap:05_design}), and the resulting implementation (\Chap{chap:06_implementation}).


% ===========================================
% ===========================================
\section{Auto-Scaling Computing Environments}
% State of the art monitoring
In recent years, container technologies have been used efficiently in complex computing environments. Dynamic scaling of containerized applications is an active area of research.
% WHich is important for this thesis
To accommodate this thesis research objective, the literature research according to auto-scaling environments was focused on two topics: Concepts of \textit{Auto-Scalers}  and auto-scaling algorithms.


\subsection{Auto-Scaler Concepts}
% A review of auto-scaling techniques for elastic application in cloud environments
\paragraph{}
Lorido-Botrán et al. \cite{Lorido2014Review} reviewed state-of-the-art literature about auto-scaling and explain auto-scaling process proposals in a cloud environment.
% SLA and costs
It is mentioned that an \textit{Auto-Scaler} is responsible to find a trade-off between meeting the Service Level Agreement (\hyperlink{abbr:sla}{SLA}) and keeping the cost of renting resources low.
% SLA
Two types of \hyperlink{abbr:sla}{SLA} exists while maintaining an acceptable trade-off: The application \hyperlink{abbr:sla}{SLA} and the resource \hyperlink{abbr:sla}{SLA}. The former is a contract between the application owner and the end-users (e.g., a certain response time). The resource \hyperlink{abbr:sla}{SLA} is agreed upon by the infrastructure provider and the application owner (e.g., 99.9\% availability).
% Auto-scaling problems
They introduced three problems \textit{Auto-Scaler} faces while scaling an environment and meeting the \hyperlink{abbr:sla}{SLA}:
% The problems
\begin{enumerate}
%1
\item Under-provisioning:
An application is under-provisioned if it needs more resources to process the incoming workload.
To make resources available and return the application to its normal state may take some time, which causes \hyperlink{abbr:sla}{SLA} violations.
% 2
\item Over-provisioning:
Applications with more resources available than needed will lead to unnecessary costs for the client.
% 3
\item Oscillation:
If scaling-actions are executed too quickly before the impact is available, a combination of over-provisioned and under-provisioned applications can occur.
A cooldown period after a scaling-action allows preventing oscillation.
\end{enumerate}
% MAPE
To prevent the mentioned problems from occurring, the authors introduced and explained the \hyperlink{abbr:mape}{MAPE} architecture (cf. \Sec{sec:02_ac}).
% MAPE phases
\hyperlink{abbr:mape}{MAPE} consists of four different phases: Monitor, Analyse, Plan, and Execute.
% Only AP
There exist \textit{Auto-Scaler} proposals which only focus on the Analyse, and Planning phase architecture of the \hyperlink{abbr:mape}{MAPE} architecture.
% Analysing and planning techniques
Several techniques for the Analyse phase are being introduced: Queuing theory and time-series analysis.
As well as for the planning phase: Threshold-based rules, reinforcement learning, and control theory.
% Active proactive grouping
There exist \textit{Auto-Scaler} which uses techniques to predict the future state of the environment (e.g., reinforcement learning). These are called proactive \textit{Auto-Scalers}.
Reactive \textit{Auto-Scalers} use techniques to respond to the environment's current status (e.g., threshold-based rules).


% Application deployment using containers with auto-scaling for microservices in cloud environment
\paragraph{}
Srirama et al. \cite{Srirama2020AppDeplyCont} designed a heuristic-based auto-scaling strategy for container-based microservices in a cloud environment. The purpose of the auto-scaling strategy is to balance the overall resource utilization across microservices in the environment.
% Results
The proposed auto-scaling strategy performed better than state-of-the-art algorithms in processing time, processing cost, and resource utilization. The processing cost of microservices was reduced by 12-20\%. The \hyperlink{abbr:cpu}{CPU} and memory utilization of cloud-servers were maximized by 9-15\% and 10-18\%, respectively.


% Comparison of auto-scaling techniques for cloud environments
\paragraph{}
Lorido-Botrán et al.  \cite{Botran2013AutoScalingComp} compared different representative auto-scaling techniques in a simulation in terms of cost and \hyperlink{abbr:sla}{SLA} violations. They compared load balancing with static threshold-based rules, reactive and proactive techniques based on \hyperlink{abbr:cpu}{CPU} load.
Load balancing is based on static rules defining the upper and lower thresholds of a specific load (e.g., \textit{if CPU > 80\% then scale-out; if CPU < 20\% then scale-in}). The difficulty of this technique is to set the ideal rules. False rules can lead to bad performance. Proactive techniques try to predict the future values of performance metrics based on historical data. Reactive techniques are based on control theory to automate the system management. To overcome the difficulties of static thresholds, the authors proposed a new auto-scaling technique using rules with dynamic thresholds. The results showed that for auto-scaling techniques to scale well, it highly depends on parameter tuning. The best result was achieved with proactive results with a minimum threshold of 20\% and a maximum of 60\%.


\subsection{Auto-Scaling Algorithms}
% Delivering elastic ...
\paragraph{}Barna et al. \cite{Barna2017ElasticContainerApps} proposed an autonomic scaling architecture approach for containerized microservices. Their approach focused on creating an autonomic management system, following the autonomic computing concept \cite{Kephart2003VisionComputing}, using a self-tuning performance model. The demonstrated architecture frequently monitors the environment and gathers performance metrics from components. It can analyze the data and dynamically scale components. In addition, to determine if a scaling action is needed, they proposed the \textit{Scaling Heat Algorithm}. The Scaling Heat algorithm is used to prevent unnecessary scaling actions, which can throw the environment temporarily off.
% Is ised in this thesis
The Scaling Heat algorithm will be used for decision making in this thesis and is explained in detail in \Sec{sec:04_background_scaling-heat}.


% Auto-scaling of Containers: the Impact of Relative and Absolute Metrics
\paragraph{}Casalicchio et al. \cite{Casalicchio2017AutoScaleCont} focused on the difference of absolute and relative metrics for container-based auto-scaling algorithms. They analysed the mechanism of the \textit{Kubernetes Horizontal Pod Auto-Scaling} (\hyperlink{abbr:khpa}{KHPA}) algorithm and proposed a new auto-scaling algorithm based on \hyperlink{abbr:khpa}{KHPA} using absolute metrics called \textit{KHPA-A}. The results showed that KHPA-A could reduce response time between 0.5x and 0.66x compared to \hyperlink{abbr:khpa}{KHPA}. In addition, their work proposed an architecture using cAdvisor for collecting container performance metrics, Prometheus for monitoring, alerting, storing time-series data, and Grafana for visualizing metrics. 
% KHPA
Absolute metrics are more appropriate when it comes to efficient resource allocation. Therefore, the KHPA-A algorithm is more efficient in vertical scaling of resources.
% This thesis
In this thesis, the focus for scaling strategies is based on the horizontal scaling approach. Therefore, the \hyperlink{abbr:khpa}{KHPA} algorithm will be used throughout this thesis and is explained in detail in \Sec{sec:04_background_khpa}.


% ===========================================
% ===========================================
\section{GPU accelerated Apache Spark Cluster}
% Short intro
This thesis aims to enable \hyperlink{abbr:gpu}{GPU} acceleration for Apache Spark.
% Research
In research, many solutions have been proposed which try to solve the problem in similar ways.
% Following
In the following, three different approaches are introduced.


% HeteroSpark
\paragraph{}
Li et al. \cite{Li2015HeteroSpark} developed a middleware framework called \textit{HeteroSpark} to enable \hyperlink{abbr:gpu}{GPU} acceleration on Apache Spark worker nodes. HeteroSpark listens for function calls in Apache Spark applications and invokes the \hyperlink{abbr:gpu}{GPU} kernel for acceleration. For communication between \hyperlink{abbr:cpu}{CPU} and \hyperlink{abbr:gpu}{GPU}, HeteroSpark implements a \hyperlink{abbr:cpu}{CPU}-\hyperlink{abbr:gpu}{GPU} communication layer for each worker node using the Java Remote Method Invocation (\hyperlink{abbr:rmi}{RMI}) \hyperlink{abbr:api}{API}. To execute operations on the \hyperlink{abbr:gpu}{GPU}, the \hyperlink{abbr:cpu}{CPU} Java Virtual Machine (\hyperlink{abbr:jvm}{JVM}) will send serialized data to the \hyperlink{abbr:gpu}{GPU} JVM using the \hyperlink{abbr:rmi}{RMI} communication interface. The \hyperlink{abbr:gpu}{GPU} \hyperlink{abbr:jvm}{JVM} will deserialize the received data for execution.
% Design
The design provides a plug-n-play approach and an \hyperlink{abbr:api}{API} for the user to call functions with \hyperlink{abbr:gpu}{GPU} support.
% Results
Overall, HeteroSpark is able to achieve an 18x speed-up for various Machine Learning applications running on Apache Spark.


% HetSpark
\paragraph{}
Klodjan et al. \cite{Klodjan2018HetSpark} introduced \textit{HetSpark}, a modification of Apache Spark.
% The goal
HetSpark extends Apache Spark with two executors, a \hyperlink{abbr:gpu}{GPU} accelerated executor and a commodity class. 
% accelerated
The accelerated executor uses VineTalk\cite{Mavridis2017VineTalk} for \hyperlink{abbr:gpu}{GPU} acceleration.
% VineTalk
VineTalk contributes as a transport layer between the application and accelerator devices (\hyperlink{abbr:cpu}{CPU} or \hyperlink{abbr:gpu}{GPU}).
% Byte code analysis
To detect suitable tasks for \hyperlink{abbr:gpu}{GPU} acceleration, HetSpark uses the ASM\footnote{ASM - \url{https://asm.ow2.io/} (Accessed: 2021-01-11)} framework to analyse the byte code of Java binaries.
% Results
The authors observed that for compute-intensive tasks, \hyperlink{abbr:gpu}{GPU} accelerated executers are preferable, while for linear tasks, \hyperlink{abbr:cpu}{CPU}-only accelerators should be used.


% Spark-GPU
\paragraph{}
% SHort intro
Yuan et al. \cite{Yuan2016SparkGPU} proposed \textit{SparkGPU} a \hyperlink{abbr:cpu}{CPU}-\hyperlink{abbr:gpu}{GPU} hybrid system built on top of Apache Spark.
% The goal
The goal of SparkGPU is to utilize \hyperlink{abbr:gpu}{GPUs} to achieve high performance and throughput.
% Problems trying to solve
SparkGPU tries to solve the following problems statements:
% The problems
\begin{enumerate}
% Iterator
\item The iterator model of Apache Spark executes one element at a time.
This approach does not match the \hyperlink{abbr:gpu}{GPU} architecture and underutilizes \hyperlink{abbr:gpu}{GPU} resources.

% JVM
\item Apache Spark runs on the \hyperlink{abbr:jvm}{JVM} and therefore stores its data on the heap memory.
\hyperlink{abbr:gpu}{GPU} programs are usually implemented with \hyperlink{abbr:gpu}{GPU} programming models like Compute Unified Device Architecture (\hyperlink{abbr:cuda}{CUDA}), which cannot access data on the heap.
Therefore, data must be copied between the heap and native memory frequently. These copy operations are expensive.

% Cluster manager
\item Existing cluster manager of Apache Spark manage \hyperlink{abbr:gpu}{GPUs} in a coarse-grained fashion.
This can lead to crashes because of insufficient memory when multiple programs run on the \hyperlink{abbr:gpu}{GPU} concurrently.
\end{enumerate}
% Extension
To solve the mentioned problem statements, SparkGPU extends Apache Spark in the following ways:
\begin{itemize}
\item Enable block processing on \hyperlink{abbr:gpu}{GPUs} by extending the Apache Sparks iterator model. Therefore Apache Spark can better utilize \hyperlink{abbr:gpu}{GPUs} to accelerate application performance.

\item To offload queries to the \hyperlink{abbr:gpu}{GPU}, SparkGPU extends the query optimizer. The query optimizer will create query plans with both \hyperlink{abbr:cpu}{CPU} and \hyperlink{abbr:gpu}{GPU} operators.

\item To manage \hyperlink{abbr:gpu}{GPUs} efficiently, SparkGPU extends the cluster manager and the task scheduler.
\end{itemize}
% Intro GPU-RDD
To extend the programming \hyperlink{abbr:api}{API}, SparkGPU provides a new Resilient Distributed Dataset (\hyperlink{abbr:rdd}{RDD}) type called GPU-RDD.
% OPtimized for GPU
A GPU-RDD is optimized to utilize the \hyperlink{abbr:gpu}{GPU}.
SparkGPU utilizes native memory on the \hyperlink{abbr:gpu}{GPU} instead of the Java heap to buffer data in GPU-RDDs.
% Operations
Operations performed on a GPU-RDD can be performed on the \hyperlink{abbr:gpu}{GPU}. Several built-in operators on the GPU-RDD are provided, which support data-parallelism.


% The query optimizer
SparkGPU can execute Structured Query Language (\hyperlink{abbr:sql}{SQL}) queries on both \hyperlink{abbr:cpu}{CPU} and \hyperlink{abbr:gpu}{GPU}.
By adding a set of \hyperlink{abbr:gpu}{GPU} rules and strategies, SparkGPU extends the query optimizer to find the best execution plan for \hyperlink{abbr:gpu}{GPU} scheduling.

% Cluster manager
To manage \hyperlink{abbr:gpu}{GPU} memory on shared \hyperlink{abbr:gpu}{GPUs}, SparkGPU provides a user-level \hyperlink{abbr:gpu}{GPU}-management library.
% Mem contention
When memory contention happens, the library will ensure that SparkGPU will stop scheduling new tasks to the Apache Spark cluster.

% End
SparkGPU accomplished to improve the performance of machine learning algorithms up to 16.13x and \hyperlink{abbr:sql}{SQL} query execution performance up to 4.83x.


% ===========================================
% ===========================================
\section{Implementation of an Automated Deployment Pipeline}
Implementing an automated deployment pipeline is a more applied topic and well described in many literatures. In this sections the primary literature used throughout the implementation of this thesis work is being introduced.


% continuous delivery
\paragraph{}The conceptual design and implementation of an automated deployment pipeline in this thesis were mostly inspired by the proposed solution of the book \textit{Continuous Delivery: Reliable Software Releases through Build, Test, and Deployment Automation} by Humble et al. \cite{Farley2010CI}.
% Short intro about the content of this book
The authors explain the theoretical idea behind an automated deployment pipeline and explaining an example implementation.
% Implementation
The proposed implementation covers the software lifecycle from compiling source code to delivering the software to a production environment.
The commit stage, which covers the build and test part of the software, can be applied in parts for this thesis work.

\chapter{Technical Background}
\label{chap:04_background}

This chapter provides information about the technologies and algorithms used for this thesis. It explains the fundamental concepts of Docker, Apache Spark, RAPIDS accelerator plugin, Prometheus, cAdvisor, and GitLab CI/CD. Additionally, the Scaling Heat, and Kubernetes Horizontal Pod Autoscaler algorithms are introduced.

% ===========================================
% ===========================================
% Docker
% ===========================================
% ===========================================
\section{Docker}
\label{sec:04_docker}
% Short docker intro
Docker is an open-source platform that enables the containerization of applications. Containerization is a technology to package, ship, and run applications and their environment in individual containers.
% Make clear
Docker is not a container technology itself. It hides the complexity of working with container technologies directly and instead provides an abstraction and the tools to work with containers \cite{Nickoloff2019Docker, Bullington2020Docker, Potdar2020Docker}.


\subsection{Docker Architecture}
\label{subsec:04_docker_architecture}
% Architecture figure
\begin{figure}[h]
\centering
\includegraphics[scale=1]{images/04_technical_background/docker/docker_architecture}
\caption{Docker architecture - Source: Authors own model, based on \cite{Docker2020Docs}.}
\label{fig:04_docker_architecture_architecture}
\end{figure}

% Explain figure
\Fig{fig:04_docker_architecture_architecture} illustrates the client-server architecture of Docker. It consists of a Docker Client, the Docker Deamon, and a Docker Registry.

\begin{itemize}
% Docker client
\item Docker Client: The Docker client is an interface for the user to send commands to the Docker daemon \cite{Docker2020Docs}.

% Docker deamon
\item Docker Deamon: The Docker daemon manages all containers running on the host system and handles container resources, networks and, volumes \cite{Bullington2020Docker}.

% Docker registry
\item Docker Registry: A Docker registry stores images. Images can be pushed to or pulled from public or private registries to build a container \cite{Docker2020Docs}.
\end{itemize}


\subsection{Docker Image}
\label{subsec:04_docker_image}
\paragraph{}
% What is an image
An Image is a snapshot of the environment to run an application in a Docker container. The environment consists of all files, libraries, and configurations that are needed for the application to run correctly \cite{Docker2020Docs, Nickoloff2019Docker}.
% How to create an image
Images can be created from existing containers or from executing a build script called \textit{Dockerfile} \cite{Nickoloff2019Docker}.


\paragraph{}
Images can be built automatically by executing a build script called Dockerfile, a text document that contains specific build instructions. Instructions are commands which are executed in order to assemble an image.
% Parent
A Dockerfile must begin with a \texttt{FROM} instruction which defines the parent image from which the image is built \cite{Docker2020Docs}.


% Example
\Lst{lst:04_docker_image_dockerfile} provides a basic example of a Dockerfile.
% Parent
In the example, \textit{ubuntu} is defined as the parent image at line 1.
% apt
Next, from 3 to 4, the Ubuntu package list is updated, and the \textit{nginx} package is installed using the \texttt{apt-get} command.
% Expose
Finally, at line 6, the Docker image is instructed to listen to port 80 at runtime with the \texttt{EXPOSE} instruction.

% Dockerfile example
\begin{lstlisting}[label=lst:04_docker_image_dockerfile, caption=Basic example of a Dockerfile]
FROM ubuntu

RUN apt-get update -qy && \
    apt-get install -y nginx

EXPOSE 80
\end{lstlisting}


\subsection{Docker Container}
% What are containers
A container is an execution environment running on the host-system kernel.

% Container image
\begin{figure}[h]
\centering
\includegraphics[scale=1]{images/04_technical_background/docker/container_structure}
\caption{Docker basic container structure - Source: Authors own model, based on \cite{Bullington2020Docker}.}
\label{fig:04_docker_container_container-structure}
\end{figure}

% Container advantages
The advantage of a container is its lightweight nature. As illustrated in \Fig{fig:04_docker_container_container-structure}, containers take advantage of \hyperlink{abbr:os}{OS}-level virtualization instead of hardware-virtualization without the need for a hypervisor \cite{Docker2020Docs, Nickoloff2019Docker}. Containers share the resources of the host system instead of using reserved resources \cite{Bullington2020Docker}. Multiple containers can run on the host-system kernel and are by default isolated from each other \cite{Docker2020Docs}.
% Docker container
In Docker, a container is a runnable unit of an image used to distribute and test applications. A container can be configured to expose specific resources to the host system, e.g., network ports \cite{Bullington2020Docker}.


\subsection{Docker Swarm Mode}
\label{subsec:04_docker_swarm}
% What is Docker swarm
Docker Swarm mode is the native cluster orchestration and management tool embedded in the Docker Engine.
% What is a swarm
In Docker Swarm mode, a cluster of multiple nodes is called a swarm. All nodes run in Swarm mode and act as managers or workers.
% Whats the purpose
In a swarm, multiple services can be deployed. The manager node is responsible for maintaining the desired state of a service \cite{Docker2020Docs}.


% 
Many properties of Docker Swarm mode contribute the fact that it is an ideal tool to create a self-healing and self-adapting environment:
\begin{itemize}
\item Desired state: The manager node monitors each service's state in the swarm and adapts the environment to maintain the desired state\cite{Docker2020Docs}.

\item Cluster management and orchestration: Docker Swarm mode is integrated within the Docker Engine. A swarm can be created and managed using the Docker command-line interface (\hyperlink{abbr:cli}{CLI}) \cite{Docker2020Docs}.

\item Service model: The Docker Engine allows to define the desired state of a service. The manager node maintains the desired state of all services in the swarm \cite{Docker2020Docs}.

\item Scaling: The number of replicas is defined for each service. The manager node automatically adapts the number of replicas for a service to keep the desired state \cite{Docker2020Docs}.

\item Multi-host networking: A swarm runs all services in an overlay network, and new services will automatically be added to it \cite{Docker2020Docs}.
\end{itemize}


\subsubsection{Services and Tasks}
% What is a service
A service defines the desired state of a task. The state is defined by the number of replicas of a service and the configuration for the Docker container, e.g., Docker image, resources, network, and more \cite{Docker2020Docs}.


%What is a node
A task is a running Docker container. The task is defined by the corresponding service and is managed by the manager node. A task can be executed on worker and manager nodes \cite{Docker2020Docs}.


\subsubsection{Nodes}
% What is a node
A Docker engine participating in the swarm is called a node.
% What about the node role
Nodes can act as manager nodes, worker nodes, or both \cite{Docker2020Docs}.


% What about the manager
The manager node is responsible for cluster orchestration and management. It maintains the desired state of all services and tasks in the swarm. In addition, the manager node dispatches tasks to worker nodes when service definitions are submitted to the manager node \cite{Docker2020Docs}.


% What about worker
Worker nodes are responsible for executing the tasks received by the manager node. While performing tasks, the worker node notifies the manager node about each task state \cite{Docker2020Docs}.


\subsubsection{Stack}
%
A stack is a named collection of first-class Docker resources (e.g., services, volumes, networks) that can be deployed to a swarm.
%
The docker-compose file format allows describing a stack in a single \hyperlink{abbr:yaml}{YAML}\footnote{YAML Ain't Markup Language - \url{https://yaml.org/} (Accessed: 2021-02-18)} file. With the \texttt{docker stack} command, a docker-compose file can be used to create a stack and deploy it to a swarm.
% Example
\Lst{lst:04_docker_swarm_stack_compose} provides an example of a stack described in the docker-compose file format.
% Explain
It defines two services called \textit{prometheus} and \textit{cadvisor}. Both services are attached to the \textit{computing-net} network.
\begin{lstlisting}[label=lst:04_docker_swarm_stack_compose, caption=Basic example of stack describes in the docker-compose file format]
version: "3.7"

networks:
    computing-net:
        name: computing_net

services:
    prometheus:
        image: prom/prometheus:latest
        networks:
            - computing-net
            
    cadvisor:
        image: google/cadvisor:latest
        networks:
            - computing-net        
\end{lstlisting}
% Deploy
\Lst{lst:04_docker_swarm_stack_deploy} shows the usage of the \texttt{docker stack} command. It deploys a stack with the name \textit{computing-stack} using a docker-compose file.
\begin{lstlisting}[label=lst:04_docker_swarm_stack_deploy, caption=Usage of the docker stack command to deploy a stack, language=bash, numbers=none]
$ docker stack deploy -c computing-stack.yml computing-stack     
\end{lstlisting}


% ===========================================
% ===========================================
% Apache Spark
% ===========================================
% ===========================================
\section{Apache Spark}
\label{sec:04_spark}
Apache Spark is an open-source computing framework for parallel data processing on a large computer cluster. Apache Spark manages the available resources and distributes computation tasks across a cluster to perform big-data processing operations at a large scale \cite{Chambers2018Spark}. Before Apache Spark was developed, Hadoop MapReduce \cite{Dean2010MapReduce} was the framework of choice for parallel operations on a computer cluster \cite{Zaharia2010Spark}. However, Spark accomplished to outperform Hadoop by 10x for iterative Machine Learning \cite{Zaharia2010Spark} and is now one of the most popular frameworks for distributed computing. It is implemented in Scala\footnote{The Scala programming language - \url{https://www.scala-lang.org/} (Accessed: 2020-12-18)}, a \hyperlink{abbr:jvm}{JVM}-based programming language. It provides a programming interface for Scala, Java\footnote{Java Software - \url{https://www.oracle.com/java/} (Accessed: 2020-12-18)}, Python\footnote{Python programming language - \url{https://www.python.org/} (Accessed: 2020-12-18)}, and R\footnote{The R Project for Statistical Computing - \url{https://www.r-project.org/} (Accessed: 2020-12-18)}. Additionally, Apache Spark includes an interactive \hyperlink{abbr:sql}{SQL} shell and libraries to implement \hyperlink{abbr:ml}{ML} and streaming applications \cite{Chambers2018Spark}.


\subsection{Spark Programming Model}
\label{subsec:04_spark_pr-model}
% Explain RDDs
Apache Spark provides Resilient Distributed Datasets (\hyperlink{abbr:rdd}{RDDs}) as the main abstraction for parallel operations \cite{Zaharia2010Spark}. Core types of the higher-level structured \hyperlink{abbr:api}{API} are built on top of \hyperlink{abbr:rdd}{RDDs}) and are automatically optimized by the Catalyst optimizer to run operations quickly and efficient \cite{Chambers2018Spark, Hien2018Spark}.


\subsubsection{Apache Spark Structured API}
% What is the structured API
Apache Spark provides high level structured \hyperlink{abbr:api}{APIs} for manipulating all kinds of data. The three distributed core types are Datasets, DataFrames, and \hyperlink{abbr:sql}{SQL}) Tables and Views \cite{Chambers2018Spark}.
% About DataFrames, Datasets and SQL
Datasets and DataFrames are immutable, lazy evaluated collections that provide execution plans for operations \cite{Chambers2018Spark}. \hyperlink{abbr:sql}{SQL} Tables and Views work the same way as DataFrames, except that \hyperlink{abbr:sql}{SQL} is used as the interface instead of using the DataFrame programming interface \cite{Chambers2018Spark}.
% Differences between Datasets ad Dataframes
Datasets use \hyperlink{abbr:jvm}{JVM} types and are therefore only available for \hyperlink{abbr:jvm}{JVM} based languages. DataFrames are Datasets of type \texttt{Row}, which is the optimized format for computations \cite{Chambers2018Spark}.


\subsubsection{Resilient Distributed Datasets}
% What are RDDs
Resilient Distributed Datasets are fault-tolerant, parallel data structures to enable data sharing across cluster applications \cite{Zaharia2012RDDs}. They express different cluster programming models like MapReduce, \hyperlink{abbr:sql}{SQL}, and batched stream processing \cite{Zaharia2012RDDs}. \hyperlink{abbr:rdd}{RDDs} have been implemented in Apache Spark and serve as the Spark structured \hyperlink{abbr:api}{API's} underlying data structure \cite{Zaharia2012RDDs}.
% How can RDDs be used in applications
\hyperlink{abbr:rdd}{RDDs} are an immutable, partitioned collection of records. They can only be initiated through transformations (e.g., map, filter) on data or other \hyperlink{abbr:rdd}{RDDs}.
% What are the advantages of RDDs
An advantage of \hyperlink{abbr:rdd}{RDDs} is that they can be recovered through lineage. Lost partitions of an \hyperlink{abbr:rdd}{RDD} can be recomputed from other \hyperlink{abbr:rdd}{RDDs} in parallel on different nodes \cite{Zaharia2012RDDs}. 
% Better use structured API
\hyperlink{abbr:rdd}{RDDs} are lower-level \hyperlink{abbr:api}{APIs} and should only be used in applications if custom data partitioning is needed \cite{Chambers2018Spark}. It is recommended to use Sparks structured \hyperlink{abbr:api}{API} objects instead. Optimizations for \hyperlink{abbr:rdd}{RDDs} have to be implemented manually, while Apache Spark automatically optimizes the execution for structured \hyperlink{abbr:api}{API} operations \cite{Chambers2018Spark}.


\subsubsection{Apache Spark Catalyst}
\label{subsubsec:04_spark_pr-model_catalyst}
% What is the Catalyst optimizer
Apache Spark also provides a query optimizer engine called Apache Spark Catalyst. \Fig{fig:04_spark_pr-model_catalyst_process} illustrates how Spark Catalyst automatically optimizes Apache Spark applications to run quickly and efficient.
% Short How
Before executing the user code, the Catalyst optimizer translates the data-processing logic into a logical plan and optimizes it using heuristics such as rules \cite{Hien2018Spark}. After that, the Catalyst optimizer converts the logical plan into a physical plan to create executable code \cite{Hien2018Spark}.


% What are logical plans
Logical plans are created from a DataFrame or a \hyperlink{abbr:sql}{SQL} query and represent the data-processing logic as a tree of operators and expressions where the Catalyst optimizer can apply sets of rule-based and cost-based optimizations \cite{Hien2018Spark}.

% What are physical plans
From the logical plan, the Catalyst optimizer creates more physical plans consisting of \hyperlink{abbr:rdd}{RDD} operations \cite{Chambers2018Spark}. The cheapest physical plan will be generated into Java bytecode for execution across the cluster \cite{Hien2018Spark}.

% Spark Catalyst figure
\begin{figure}[h]
\centering
\includegraphics[scale=1]{images/04_technical_background/spark/catalyst_optimization_process}
\caption{Optimization process of the Spark Catalyst - Source: Authors own model, based on \cite{Hien2018Spark}.}
\label{fig:04_spark_pr-model_catalyst_process}
\end{figure}


\subsection{Application Architecture}
\label{subsec:04_spark_architecture}

\begin{figure}[h]
\centering
\includegraphics[scale=1]{images/04_technical_background/spark/spark_cluster_architecture}
\caption{Overview of a Spark cluster architecture - Source: Authors own model, based on \cite{Apache2020Spark}.}
\label{fig:spark_cluster_overview}
\end{figure}

% Explain image
\Fig{fig:spark_cluster_overview} illustrates the main architecture of an Apache Spark cluster running an application.
% Master-Worker
The architecture follows the master-worker model.
% The processes
Each running application has one driver process (master) and multiple executor processes (worker) exclusively assigned by the cluster manager.
% Cluster manager
Furthermore, the cluster manager decides on which nodes the processes will be executed \cite{Hien2018Spark}.


\subsubsection{Driver Process}
\label{subsubsec:04_spark_architecture_driver}
The driver process is a \hyperlink{abbr:jvm}{JVM} process running on a physical machine and is responsible for maintaining the execution of an Apache Spark application \cite{Chambers2018Spark}. It coordinates the application tasks and schedules them on available executors \cite{Hien2018Spark}. 
The driver interacts with the cluster manager to launch executors and allocate hardware resources \cite{Chambers2018Spark, Hien2018Spark}.


\subsubsection{Executor Process}
% JVM
The executor process is a \hyperlink{abbr:jvm}{JVM} process that runs through the whole duration of an application \cite{Hien2018Spark, Apache2020Spark}.
% Tasks
It is responsible for performing all tasks (units of work) assigned by the driver process \cite{Chambers2018Spark}.
% Report
After the executor process finishes, it reports back to the driver process \cite{Chambers2018Spark}.
% CPU
Each task is performed on a separate \hyperlink{abbr:cpu}{CPU} core to enable parallel processing \cite{Hien2018Spark}.


\subsubsection{Cluster Manager}
\label{subsubsec:04_spark_architecture_manager}
The cluster manager is an external service that orchestrates the work between available machines in the cluster \cite{Hien2018Spark, Apache2020Spark}.
% Worker and master nodes
It decides on which nodes the driver- and executor processes are launched in the cluster.
% Resources
Additionally, the cluster manager maintains each node's resources in the cluster \cite{Hien2018Spark, Chambers2018Spark}.
% Different cluster types
Apache Spark supports different services that can run as cluster manager, such as Standalone mode (introduced in \Sec{subsec:04_spark_standalone}), Apache Mesos\footnote{Apache Mesos - \url{https://mesos.apache.org/} (Accessed: 2021-01-02)}, Hadoop YARN\cite{Murthy2013Yarn}, and Kubernetes\footnote{Kubernetes - \url{https://kubernetes.io/} (Accessed: 2021-01-02)} \cite{Apache2020Spark}.
% Explain different cluster deploy modes
The cluster manager provides three different deploy modes for acquiring resources in the cluster.
\begin{itemize}
% Cluster mode
\item Cluster mode:
To run an application in cluster mode, the user must submit a precompiled Java ARchive (\hyperlink{abbr:jar}{JAR}), python script, or R script to the cluster manager \cite{Chambers2018Spark}. After that, the cluster manager starts the driver process and executor processes exclusively for the Apache Spark application on machines inside the cluster \cite{Chambers2018Spark, Hien2018Spark}.

% Client mode
\item Client mode:
The difference between the client mode and the cluster mode is that, the driver process runs on the client machine outside of the Spark cluster \cite{Chambers2018Spark}.

% Local mode
\item Local mode:
The local mode starts an Apache Spark application on a single computer \cite{Chambers2018Spark}. It is important to mention that the local mode is not recommended to use in production. Instead, it should be used for testing Apache Spark applications \cite{Chambers2018Spark}.
\end{itemize}


\subsection{Standalone Cluster Deployment}
\label{subsec:04_spark_standalone}
% What is a standalone cluster
The standalone mode is a basic cluster-manager build specifically for Apache Spark. It is developed only to run Apache Spark but supports workloads at a large scale \cite{Chambers2018Spark}.
% Why standalone cluster mode, advantages


\subsubsection{Starting Master and Worker Nodes}
\paragraph{}
Apache Spark provides executable launch scripts to start master and worker nodes in standalone mode.
% Where to find them
Inside the Apache Spark installation directory, the executables can be found at \texttt{sbin/start-master.sh} to start a master node and at \texttt{/start-slave.sh} to start a worker node.
% parameter for worker script
The worker launch executable requires the master node Uniform Resource Identifier (\hyperlink{abbr:uri}{URI}) as a parameter \cite{Apache2020Spark}. 


% Example master
\begin{lstlisting}[label=lst:04_spark_standalone_launch-master, caption=Usage of master launch script, language=bash, numbers=none]
$ ./sbin/start-master.sh
\end{lstlisting}


% Example worker
\begin{lstlisting}[label=lst:04_spark_standalone_launch-worker, caption=Usage of worker launch script, language=bash, numbers=none]
$ ./sbin/start-slave.sh spark://spark-master:7077
\end{lstlisting}


% Introduce examples
\paragraph{}\Lst{lst:04_spark_standalone_launch-master} and \Lst{lst:04_spark_standalone_launch-worker} provide an example of using both executables to start a master and a worker node. The \hyperlink{abbr:uri}{URI} \texttt{spark://spark-master:7077} in \Lst{lst:04_spark_standalone_launch-worker} is an example of a master node \hyperlink{abbr:uri}{URI}. The master node launch script will print out the master \hyperlink{abbr:uri}{URI} after being executed successfully \cite{Apache2020Spark}.


\subsubsection{Resource Allocation}
\label{subsubsec:04_spark_standalone_res-alloc}
In standalone mode, workers need a set of resources configured. Therefore, a worker can assign resources to executors.
% Discovery script
To specify how a worker discovers resources, a discovery script has to be available \cite{Apache2020Spark}.


\subsubsection{Submitting Applications with spark-submit}
\label{subsubsec:04_spark_standalone_submit}
% Which executable
\paragraph{}
To submit an application to a standalone cluster, Apache Spark provides the \texttt{spark-submit} executable. The executable file is available at \texttt{bin/spark-submit} in the installation folder of Apache Spark.
% Cluster mode
In cluster mode, the driver of an Apache Spark application (see \Sec{subsubsec:04_spark_architecture_driver}) is launched from one of the worker processes inside the cluster.
% When does it finish
The submit process will finish after it has submitted the application. It does not wait for the submitted application to finish \cite{Apache2020Spark}.


% Example
\begin{lstlisting}[label=lst:04_spark_standalone_submit_example, caption=Example usage of the spark-submit executable, language=bash, numbers=none]
$ bin/spark-submit --master spark://spark-master:7077 application.py
\end{lstlisting}


% How to use it
\paragraph{}
\Lst{lst:04_spark_standalone_submit_example} shows how the spark-submit executable can be used to submit a Python application to a standalone Apache Spark cluster.
% required parameter
Spark-submit requires the master node \hyperlink{abbr:uri}{URI} and the path to the desired Spark application file.
% Which files
With the spark-submit executable, it is possible to submit Python, Java, and R applications \cite{Apache2020Spark}.


% ===========================================
% ===========================================
% NVIDIA RAPIDS
% ===========================================
% ===========================================
\section{RAPIDS Accelerator for Apache Spark}
\label{sec:04_rapids}
% Explain
RAPIDS accelerator for Apache Spark is a plugin suite that aims to accelerate computing operations for Apache Spark by executing pipelines entirely on \hyperlink{abbr:gpu}{GPUs}. It is available for Apache Spark 3 \cite{SparkRapids2020Docs}.
% A little how
The plugin uses the RAPIDS AI cuDF\footnote{Open GPU Data Science - \url{https://rapids.ai/} (Accessed: 2021-01-01)} library to extend the Apache Spark programming model, introduced in \Sec{subsec:04_spark_pr-model} \cite{SparkRapids2020Docs, Mcdonald2020SparkRapids, Aguerzame2019GPUAO}.


\subsection{Extension of the Spark programming model}
\label{subsec:04_rapids_ext}
% How does it extends DataFrame and SQL
\paragraph{}The plugin suite extends the Apache Spark programming model with a new DataFrame based on Apache Arrow\footnote{Arrow. A cross-language development platform for in-memory data - \url{https://arrow.apache.org/} (Accessed: 2020-12-03)} data structures. The new DataFrame \hyperlink{abbr:api}{API} aims to accelerate loading, filtering, and manipulation operations on large datasets. In addition, the Catalyst optimizer (described in \SubSec{subsubsec:04_spark_pr-model_catalyst}) is extended to generate \hyperlink{abbr:gpu}{GPU}-aware query plans \cite{Mcdonald2020SparkRapids, Aguerzame2019GPUAO}.
% What is Apache Arrow
Apache arrow is a data platform to build high-performance applications that work with large datasets and improve analytic algorithms. A component of Apache Arrow is the Arrow Columnar Format, an in-memory data structure specification for efficient analytic operations on \hyperlink{abbr:gpu}{GPUs} and \hyperlink{abbr:cpu}{CPUs} \cite{ApacheArrow2020Docs}.


% RAPIDS query plan image
\begin{figure}[h]
\centering
\includegraphics[scale=1]{images/04_technical_background/rapids/rapids_catalyst_optimization_process}
\caption{Catalyst optimization with RAPIDS accelerator for Apache Spark - Source: Authors own model, based on \cite{Mcdonald2020SparkRapids}.}
\label{fig:04_rapids_ext_query-plan}
\end{figure}
\paragraph{}\Fig{fig:04_rapids_ext_query-plan} illustrates how the RAPIDS plugin suite extends the Catalyst optimization process illustrated previously in \Fig{fig:04_spark_pr-model_catalyst_process}.
% How does the Catalyst create GPU aware plans
The Spark Catalyst optimizer identifies operators in a query plan that are supported by the RAPIDS \hyperlink{abbr:api}{API}. To execute the query plan, these operators can be scheduled on a \hyperlink{abbr:gpu}{GPU} within the Spark cluster \cite{Mcdonald2020SparkRapids}. If the RAPIDS \hyperlink{abbr:api}{APIs} do not support operators, a physical plan for \hyperlink{abbr:cpu}{CPUs} is generated by the Catalyst optimizer to execute \hyperlink{abbr:rdd}{RDD} operations \cite{Mcdonald2020SparkRapids}.


\subsection{GPU Accelerated Machine Learning with XGBoost}
% SHort Intro
RAPIDS accelerates SparkSQL operations and operations on a DataFrame. Additionally, RAPIDS aims to accelerate the training process of machine learning models.
% Only XGBoost
Currently, RAPIDS only supports \hyperlink{abbr:gpu}{GPU}-acceleration for Extreme Gradient Boosting (\hyperlink{abbr:xgboost}{XGBoost}) in SparkML \cite{Mcdonald2020SparkRapids}.


% Whats XGBoost
\hyperlink{abbr:xgboost}{XGBoost} is a scalable, distributed gradient-boosted machine learning library. It tries to solve many data science problems, by implementing \hyperlink{abbr:ml}{ML} algorithms using the gradient boosting technique.
% XGBoost4j-Spark
With the XGBoost4j-Spark\footnote{XGBoost Documentation - \url{https://xgboost.readthedocs.io/en/latest/jvm/xgboost4j_spark_tutorial.html} (Accessed: 2021-01-28)} library, \hyperlink{abbr:xgboost}{XGBoost} integrates into the Apache SparkML library \cite{XGBoost2021Docs}. 


\subsection{Installation Requirements for Apache Spark Standalone Mode}
\label{subsec:04_rapids_req}
The RAPIDS accelerator for Apache Spark is available for a standalone mode Apache Spark cluster. To operate efficiently, the following requirements need to be installed on all Apache Spark nodes in the cluster \cite{SparkRapids2020Docs}:
% requirements
\begin{itemize}
\item Java Runtime Environment (\hyperlink{abbr:jre}{JRE})
\item NVIDIA \hyperlink{abbr:gpu}{GPU} driver
\item \hyperlink{abbr:cuda}{CUDA} Toolkit\footnote{CUDA Toolkit - \url{https://developer.nvidia.com/cuda-toolkit} (Accessed: 2021-01-01)}
\item RAPIDS accelerator for Apache Spark Java library
\item cuDF Java library that is supported by the RAPIDS accelerator Java library, and the installed \hyperlink{abbr:cuda}{CUDA} toolkit
\item XGBoost4j Java library
\item XGBoost4j-Spark Java library
\item \hyperlink{abbr:gpu}{GPU} resource discovery script
\end{itemize}


% ===========================================
% ===========================================
% PROMETHEUS
% ===========================================
% ===========================================
\section{Prometheus}
\label{sec:04_prom}
% What is Prometheus
Prometheus is an open-source monitoring and alerting system \cite{Prom2020Docs}.
% Time-series format + PromQL
To collect and store data, Prometheus supports a multi-dimensional key-value pair based data model, according to \SubSec{subsec:02_monitoring_db_multi-model}, which can be analyzed by using the PromQL query language \cite{Pandey2020Monitoring}.
% SHort PromQL
PromQL is a functional query language for selecting and aggregating time-series data in real-time \cite{Prom2020Docs}.
% Pull-based approach
Prometheus follows the pull-based approach described in \SubSec{subsec:02_monitoring_push-pull} to scrape metrics from hosts and services \cite{Bastos2019Prom}.


\subsection{Prometheus Architecture}
\label{sec:04_prom_arch}
% Architecture image
\begin{figure}[h]
\centering
\includegraphics[scale=0.8]{images/04_technical_background/prometheus/prometheus_architecture}
\caption{Prometheus high-level architecture - Source: Authors own model, based on \cite{Prom2020Docs, Brazil2018Prom}.}
\label{fig:prom_architecture}
\end{figure}

% High level architecture
\Fig{fig:prom_architecture} illustrates the high-level architecture of Prometheus.
% Description of components
The Prometheus ecosystem provides multiple components, which can be optional, depending on the environment's monitoring needs \cite{Bastos2019Prom}. The main components of a Prometheus system are Prometheus Server, Alertmanager, service discovery, exporters, Push Gateway, and data visualization \cite{Prom2020Docs}.

\subsubsection{Prometheus Server}
% SHort intro
The Prometheus server is the main component of a Prometheus system. It is responsible for collecting metrics as time-series data from targets and stores the collected data in the built-in \hyperlink{abbr:tsdb}{TSDB} \cite{Bastos2019Prom}. Prometheus uses the concept of scraping to collect metrics from a target. A target host has to expose an endpoint to make metrics available in the Prometheus data format \cite{Pandey2020Monitoring}. Additionally, the Prometheus server triggers alerts to the Alertmanager if a configured condition becomes true \cite{Prom2020Docs}.
% The figure
The core components of the Prometheus server, as illustrated in \Fig{fig:prom_architecture}, are the following:

\begin{itemize}
% Service discovery
\item Service Discovery:
As being mentioned before, Prometheus follows a pull-based approach to fetch metrics from a target. To know about all targets, Prometheus needs a list of the corresponding hosts. The service discovery manages the complexity of maintaining a list of hosts manually in a changing infrastructure \cite{Bastos2019Prom}. Therefore, Prometheus is able to notice targets that are not responding \cite{Brazil2018Prom}.

% Retrieval
\item Retrieval:
Prometheus sends an Hypertext Transfer Protocol (\hyperlink{abbr:http}{HTTP}) request to each target to scrape metrics. The request is sent each interval, which can be set in the configuration \cite{Brazil2018Prom}.

% HTTP API
\item \hyperlink{abbr:http}{HTTP} \hyperlink{abbr:api}{API}:
Prometheus provides an \hyperlink{abbr:http}{HTTP} \hyperlink{abbr:api}{API}. This \hyperlink{abbr:api}{API} can be used to request raw data and evaluate PromQL queries.
A data visualization tool can use this \hyperlink{abbr:api}{API} to create visualizations of the requested metrics \cite{Brazil2018Prom}.

% Recording rules and alerts
\item Recording and alert rules:
Recording rules enable to precompute frequently needed or compute-intensive PromQL expressions. The result will be saved as a set of time-series in the local storage. This enables to query a recording rule at a much faster speed than the original PromQL expression \cite{Brazil2018Prom, Prom2020Docs}.

Alert rules define conditions based on PromQL expressions. If a condition becomes true, an alert will be sent to an external service \cite{Prom2020Docs}.

% Storage
\item Storage:
Received data is stored in a custom highly efficient format on a local on-disk time-series database \cite{Prom2020Docs}. Prometheus does not offer a solution for distributed storage across a cluster of machines \cite{Brazil2018Prom}.
\end{itemize}


\subsubsection{Optional Components}
The Prometheus ecosystem offers a set of optional components and can be activated depending on the monitoring needs.
The optional components illustrated in \Fig{fig:prom_architecture} are the following:

\begin{itemize}
% Alertmanager
\item Alertmanager:
If an alerting rule becomes true, the Prometheus server generates an alert and pushes it to the Alertmanager. The Alertmanager generates notifications from the received alerts. A notification can take multiple forms like emails or chat messages. Webhooks can be implemented to trigger custom notifications \cite{Bastos2019Prom}.

% Exporters
\item Exporters:
If an application does not support an endpoint for Prometheus, an exporter can be used to fetch metrics and make them available to the Prometheus server. An exporter is a monitoring agent running on a target host that fetches metrics from the host and exports them to the Prometheus server \cite{Pandey2020Monitoring}.

% Push Gateway
\item Push Gateway:
If a target is not designed to be scraped, metrics can be pushed against the Push Gateway \cite{Prom2020Docs}. The Push Gateway converts the data into the Prometheus data format and passes them to the Prometheus server \cite{Pandey2020Monitoring}.

% Data vis
\item Data Visualisation:
Prometheus supports various visualization tools, such as its built-in visualization or third-party tools. One of the widely used tools for this occasion is Grafana\footnote{Grafana: The open observability platform - \url{https://grafana.com/} (Accessed: 2021-01-19)}.
\end{itemize}


\subsection{Prometheus Configuration}
\label{sec:04_prom_config}
% Short intro
Prometheus scraping jobs and rules can be configured using external \hyperlink{abbr:yml}{YAML} files \cite{Prom2020Docs}.
% Explain listing
\Lst{lst:04_prom_config_example} shows a valid configuration file example.
% global section
In the \texttt{global} configuration section, default values can be set.
% scrape_configs
Targets are defined in the \texttt{scrape\_configs} section. Each target is defined as a scrape job with a unique name.  A target can be defined statically using the \texttt{static\_configs} parameter or dynamically using the available service discovery mechanisms \cite{Prom2020Docs}.
% rules
Rules have to be configured in a separate \hyperlink{abbr:yml}{YAML} file. To load rules into Prometheus, the file path has to be set in the \texttt{rule\_files} parameter.

% config file example
\begin{lstlisting}[label=lst:04_prom_config_example, caption=Prometheus configuration file example]
global:
  scrape_interval: 5s

scrape_configs:
  - job_name: cadvisor
    static_configs:
      - targets: ["cadvisor:8080"]
        labels:
          group: "cadvisor"
          
rule_files:
  - "/etc/prometheus/recording_rules.yml"
\end{lstlisting}

% rules example
\Lst{lst:04_prom_config_rule-example} shows a configuration example of a recording rule which are defined in groups. 
% Rule config
Each rule is defined by a name and a valid PromQL expression \cite{Prom2020Docs}.

% Recording rule example
\begin{lstlisting}[label=lst:04_prom_config_rule-example, caption=Prometheus rules configuration file example]
groups:
  - name: http_requests
     - record: job:http_inprogress_requests:sum
       expr: sum by (job) (http_inprogress_requests)
\end{lstlisting}


% ===========================================
% ===========================================
% CADVISOR
% ===========================================
% ===========================================
\section{cAdvisor}
% What is cAdvisor
Container Advisor (cAdvisor) is a running daemon that collects, aggregates, analyses, and exposes performance metrics from running containers.
% How to run
It has native support and is deployed as a Docker container.
% How it collects
cAdvisor collects metrics from the container daemon and Linux cgroups and exposes them in the Prometheus file format \cite{Bastos2019Prom, cadvisor2020Docs}.


% ===========================================
% ===========================================
% DCGM-EXPORTER
% ===========================================
% ===========================================
\section{DCGM-Exporter}
% Intro
DCGM-Exporter is a monitoring agent for collecting health and performance metrics from NVIDIA \hyperlink{abbr:gpu}{GPUs}. It exports collected metrics in the Prometheus format and makes them available using an \hyperlink{abbr:http}{HTTP} endpoint \footnote{DCGM-Exporter - \url{https://ngc.nvidia.com/catalog/containers/nvidia:k8s:dcgm-exporter} (Accessed: 2021-02-12)}.
% DCGM
The DCGM-Exporter is built on top of the NVIDIA Data Center GPU Manager (\hyperlink{abbr:dcgm}{DCGM}), a tool-suite for managing and monitoring NVIDIA \hyperlink{abbr:gpu}{GPUs} in a cluster \footnote{NVIDIA DCGM - \url{https://developer.nvidia.com/dcgm} (Accessed: 2021-02-12)}.
%
It can be deployed to a cluster using a Docker container. \Lst{lst:04_dcgm_deploy} provides an example of deploying the DCGM-Exporter using a Docker container. As a requirement, the NVIDIA Container Toolkit\footnote{NVIDIA Cloud Native technologies documentation - \url{https://docs.nvidia.com/datacenter/cloud-native/container-toolkit/overview.html} (Accessed: 2021-01-28)} has to be installed on the host.
\begin{lstlisting}[label=lst:04_dcgm_deploy, caption=DCGM-Exporter deployment using a Docker container, language=bash, numbers=none]
$ docker run -d --rm \
    -p 9400:9400 \
    --runtime nvidia \
    nvidia/dcgm-exporter:2.0.13-2.1.1-ubuntu18.04
\end{lstlisting}


% ===========================================
% ===========================================
% Gitlab CI/CD
% ===========================================
% ===========================================
\section{GitLab CI/CD}
\label{sec:04_background_gitlab}
GitLab CI/CD is a tool integrated into the GitLab platform that enables Continuous Integration (\hyperlink{abbr:ci}{CI}) and Continuous Delivery/Deployment (\hyperlink{abbr:cd}{CD}) for software development.
% Gitlab platform
The GitLab platform integrates many development features like Git repository management and \hyperlink{abbr:ci}{CI}/\hyperlink{abbr:cd}{CD}.
% How it works
By pushing code changes to the codebase, GitLab CI/CD executes a pipeline of scripts to automate the software development cycle's \hyperlink{abbr:ci}{CI} and \hyperlink{abbr:cd}{CD} processes.
% CI
A \hyperlink{abbr:ci}{CI} pipeline consists of scripts that build, test, and validate the updated codebase.
% CD
A \hyperlink{abbr:cd}{CD} pipeline is responsible for deploying the application for production after the \hyperlink{abbr:ci}{CI} pipeline was executed successfully.
% Summary
Adding \hyperlink{abbr:ci}{CI}/\hyperlink{abbr:cd}{CD} pipelines to the software development cycle of an application allows to catch bugs and errors early. This ensures that an application deployed to production will conform to established standards \cite{Gitlab2020Docs}.

\subsection{CI/CD Pipeline}
\label{subsec:04_background_gitlab_pipeline}
The fundamental component of GitLab CI/CD is called a pipeline. Pipelines perform based on conditions. A condition might be a push to the main branch or a specific branch of the repository \cite{Gitlab2020Docs}. A pipeline comprises of two components:

\begin{itemize}
\item Stages: A stage consists of one or multiple jobs that run in parallel. Furthermore, a stage defines how jobs will be executed. For example, a build stage only performs after a test stage has performed successfully \cite{Gitlab2020Docs}.

\item Jobs: Jobs are responsible for performing the scripts defined by administrators. The scripts define necessary actions. For example compiling the source code or performing tests \cite{Gitlab2020Docs}.
\end{itemize}


% How to configurate a pipeline (.gitlab-ci.yml)
GitLab CI/CD is configured by a \texttt{.gitlab-ci.yml} file. The file must be located in the repository root directory.
% What does this file
The configuration file will create a pipeline that performs after a push to the repository \cite{Gitlab2020Docs}.


\subsection{Example of a Basic CI/CD Pipeline Architecture}
\label{subsec:04_background_gitlab_basic-pipeline}
\begin{figure}[h]
\centering
\includegraphics[scale=0.8]{images/04_technical_background/gitlab/basic_pipeline}
\caption{Basic architecture of a GitLab CI/CD pipeline - Source: Authors own model, based on \cite{Gitlab2020Docs}.}
\label{fig:04_gitlab_pipeline_basic_arch}
\end{figure}

% Explain figure
\Fig{fig:04_gitlab_pipeline_basic_arch} illustrates an architecture of a basic \hyperlink{abbr:ci}{CI}/\hyperlink{abbr:cd}{CD} pipeline.
% Two stages
The pipeline consists of a build and a test stage, which are performed in order. The test stage will only be performed after the build stage was successful.
% Build stage
The build stage consists of a single job named \textit{build-job}, which is executed first after a change is pushed to the repository.
% Test stage
The test stage consists of two jobs called \textit{unit-test-job} and \textit{integration-test-job}. If a stage consists of more than one job, then all jobs are executed in parallel when the stage is triggered.


% Example config

% Explain example
\Lst{lst:04_gitlab_pipeline_basic_config-example} provides the configuration for the basic pipeline example.
% Shell script
Each job executes a shell script to perform the desired actions. The shell scripts have to be located in the source code repository.
\newpage
% The example
\begin{lstlisting}[label=lst:04_gitlab_pipeline_basic_config-example, caption=Example of a \texttt{.gitlab-ci.yml} configuration file]
stages:
  - build
  - test

build-job:
  stage: build
  script:
    - build_software.sh

unit-test-job:
  stage: test
  script:
    - run_unit_tests.sh
    
integration-test-job:
  stage: test
  script:
    - run_integration_tests.sh
\end{lstlisting}


\subsection{Job Execution}
\label{sec:04_background_gitlab_job}
% Runner pick up jobs
GitLab runners will perform jobs that are defined in the configuration file. 
% What is runner
A GitLab runner is a lightweight and highly scalable application that runs on a server and performs one or multiple executors. Furthermore, it performs given jobs in its environment and responds the result back to the GitLab instance.
% WHat is an executor
An executor provides the environment where jobs are executed in. GitLab runner provides multiple executors, such as the Docker executor that connects to the underlying Docker engine. In addition, the Docker executor performs a job in a separate and isolated Docker container.
% Runner config
GitLab runner can be set up for specific projects or be available to all projects on the GitLab platform \cite{Gitlab2020Docs}.


% ===========================================
% ===========================================
% Scaling heat
% ===========================================
% ===========================================
\newpage
\section{Scaling Heat}
\label{sec:04_scal-heat}
% Intro
The Scaling Heat algorithm is a decision making algorithm to determine if a scaling action is necessary.
% Authors and why
It was introduced by Barna et al. \cite{Barna2017ElasticContainerApps} to overcome issues of traditional recurrence factor algorithms \cite{Barna2017ElasticContainerApps}.


\subsection{Recurrence Factor}
% Why recurrence factor
\paragraph{}
In an autonomic computing environment, a scaling decision is made in each interval after data has been retrieved from the monitoring system (see \Sec{sec:02_ac} for the autonomic computing architecture). 
% Performance spikes
Sudden performance spikes can occur, and the decision algorithm performs unnecessary scaling actions.
% The problem
These unnecessary scaling actions can have a negative impact on the overall computing performance.
% The recurrence factor
To overcome this issue, a recurrence factor needs to be introduced to the decision making algorithm.
% What it does
With a recurrence factor ($n$), a scaling action is only performed until a performance threshold is violated $n$ times \cite{Barna2017ElasticContainerApps}.


% Problems of traditional recurrence factor algos
\paragraph{}Traditional recurrence factor algorithms require violations to occur regularly. If a performance violation of the opposite direction occurs, the algorithm can get stuck in the violation process. Then, no scaling actions are performed \cite{Barna2017ElasticContainerApps}.


\subsection{Scaling Heat Algorithm Concept}
\label{sec:04_background_scaling-heat}
% Describe the concept
\paragraph{}\Alg{alg:04_scal-heat_concept_algo} introduces the Scaling Heat algorithm.
% The base
The algorithm is based on a concept called \textit{heat}. The value of heat indicates if a scaling action of adding or removing components is necessary.
% Violations
If the given utilization of a performance metric violates the upper threshold, the heat value will increase. Violations of the lower threshold will cause a decrease, respectively.
% Recurrence
When the heat reaches the recurrence factor $n$, positive for adding and negative for removing nodes, a scaling action is executed.
% Reset
After executing a scaling action, the heat value is reset to 0 \cite{Barna2017ElasticContainerApps}.
% The algorithm
\begin{algorithm}[H]
\SetKwInput{KwInput}{Input}
\SetKwInput{KwOutput}{Output}
\SetAlgoLined
\KwInput{$utilization$ - The retrieved utilization for a performance metric}
\KwInput{$lower\_threshold$ and $upper\_threshold$ - Range limit of the performance metric}
\KwInput{$heat$ - Current heat value of a performance metric. Indicating if a scaling action is necessary}
\KwOutput{$heat$ - New heat value for the next iteration}
\uIf{$utilization \geq upper\_threshold$}{
   \tcp{cluster overload}
   \uIf{$heat < 0$}{
     \tcp{reset heat for removal}
     $heat \leftarrow 0$\;
   }
   $heat \leftarrow heat + 1$\;
   }
   \uElseIf{$utilization \leq lower\_threshold$}{
     \tcp{cluster overload}
     \uIf{$heat > 0$}{
       \tcp{reset heat for adding}
       $heat \leftarrow 0$\;
     }
     $heat \leftarrow heat - 1$\;
   }  
   \Else{
   \tcp{utilization is within threshold range}
   \tcp{move heat towards 0}
   \uIf{$heat > 0$}{
      $heat \leftarrow heat - 1$\;
   }
   \uElseIf{$heat < 0$}{
      $heat \leftarrow heat + 1$\;   
   }
  }
\uIf{$heat = n$}{
Perform a scale-out action\;
$heat \leftarrow 0$\;
}
\uElseIf{$heat = -n$}{
Perform a scale-in action\;
$heat \leftarrow 0$\;
}
\Return{heat}
\caption{Scaling Heat decision making algorithm\cite{Barna2017ElasticContainerApps}}
\label{alg:04_scal-heat_concept_algo}
\end{algorithm}


% ===========================================
% ===========================================
% KHP
% ===========================================
% ===========================================
\newpage
\section{Kubernetes Horizontal Pod Autoscaler}
\label{sec:04_background_khpa}
\paragraph{}
% Short intro
Kubernetes Horizontal Pod Autoscaler (\hyperlink{abbr:khpa}{KHPA}) is an auto-scaling algorithm used in Kubernetes, which is an orchestration tool that allows to create and deploy units called Pods. A Pod is a running process on a cluster that encapsulates an application.
% What
\hyperlink{abbr:khpa}{KHPA} scales the number of replicas per Pod based on the utilization of performance metrics.
% Interval
The algorithm is based on a control loop. Each $n$ seconds, the algorithm gathers performance metrics and computes the target number of replicas to achieve the desired utilization of a performance metric \cite{Casalicchio2017AutoScaleCont}.


% Support for multiple metrics
The algorithm computes the number of replicas for a single performance metric. If a scaling action depends on multiple performance metrics, the number of replicas has to be computed for each performance metric. The largest number of replicas is used as the target number of replicas \cite{Kubernetes2021Docs}.


\paragraph{}
\hyperlink{abbr:khpa}{KHPA} takes as input the number of active replicas for a pod ($active\_replicas$), the utilization of the performance metric of each replica ($pod\_utilization$), and the target utilization of the performance metric ($target\_utilization$).
% Intro formula
The formula to compute the target number of pods $P$ is defined by \cite{Kubernetes2021Docs}:
% The formula
\begin{equation}
P = \left \lceil active\_replicas \times \left ( \frac{\sum pod\_utilization}{target\_utilization} \right ) \right \rceil
\label{eq:04_background_khpa_equation}
\end{equation}


\chapter{Conceptual Design}
\label{sec:design}
%
\todo{Describe Chapter}


% ===========================================
% ===========================================
\section{Design Restrictions}
\label{sec:05_restrictions}
The design of the computing environment will be restricted by several points. These factors are given because of technological choices and requirements from up there.

\begin{itemize}
\item \textbf{Running on a NVIDIA DGX A100\footnote{The Universal System for AI Infrastructure - \url{https://www.nvidia.com/en-us/data-center/dgx-a100/}}:} The environment will run on a NVIDIA DGX A100 workstation. The workstation has 80 CPUs and 8 GPUs installed. For this computing environment, 2 GPUs will be available.

\item \textbf{Apache Spark for distributed computing:} Apache Spark will be used as a distributed computing framework.

\item \textbf{Python as the main programming language:} Python will be used as the main programming language for Apache Spark applications. Therefore, examples will use Python code. Configurations for the system are optimized for using Python in production.
\end{itemize}

% ===========================================
% ===========================================
\section{CI/CD}


% ===========================================
% ===========================================
\section{Identification of suitable Metrics for Scaling}
% Intro sentence
To analyze the Apache Spark cluster computing performance suitable metrics have to be defined.
% How do we compute
As mention in SECTION XY, Apache Spark distributes its workload across multiple CPU cores. In addition, one objective of this thesis is to accelerate the computing performance of the Apache Spark cluster with GPUs.
% List metrics
Therefore suitable performance metrics are CPU utilization of all Apache Spark worker and the utilization of all available GPUs.
% Those are time-based utilization
These utilization metrics will be based on the time, when the Apche Spark cluster is actively performing computations.


\subsection{CPU Performance}
% Evtl ist hier der faösche platz -> Besser eine Metrics section wo erklärt wird, wie das berechnet werden soll.
\todo{Hier auf Tabelle verweisen (Anhang) Metriken die von Prometheus + cAdvisor bereitgestellt werden}

% How is the CPU Utilization being calculated
To adapt to business needs, the CPU percentage of each Spark Worker will be calculated. Prometheus provides several metrics to calculate the CPU percentage. The CPU percentage of all Worker can be calculated as follows:

Each Apache Spark worker has a number of CPU cores to perform work. To calculate the CPU utilization of an Apache Spark worker,

\begin{equation}
SparkWorkerCPUUtilization = \dfrac{\sum CPUCoreUtilization}{NumberOfCPUCores}
\label{eq:formel}
\end{equation}

\begin{equation}
OverallCPUUtilization = \dfrac{\sum SparkWorkerCPUUtilization}{NumberOfSparkWorker}
\label{eq:formel}
\end{equation}

\subsection{GPU Performance}
The system has a fixed number of GPUs to use.
\begin{equation}
GPUUtilization = \dfrac{\sum SparkWorkerGPUUtilization}{NumberOfSparkWorker}
\label{eq:formel}
\end{equation}


% ===========================================
% ===========================================
\section{Computing environment Architecture}


\subsection{Overall}
\label{subsec:05_arch_overall}

% Figure
\begin{figure}[h]
\centering
\includegraphics[scale=0.8]{images/05_conceptual_design/cluster_architecture/overall_architecture}
\caption{Overall cluster architecture - Source: Authors own model.}
\label{fig:ca-overall_architecture}
\end{figure}

\Fig{fig:ca-overall_architecture} illustrates the overall architecture with all components of the computing environment. The two main components in the environment are an Apache Spark cluster and an autonomic manager.
% The components
The autonomic manager will be implemented according to the MAPE architecture (introduced in \SubSec{subsec:02_ac_manager}). It is responsible for monitoring and auto-scaling the Apache Spark cluster. To enable distributed computing, an Apache Spark cluster will be set up to execute machine learning applications. 
% The nodes and virtualization
The autonomic manager and the Apache Spark cluster consist each of multiple nodes. Each node of a component will run as a Docker container.
% Self healing
All containers will run as services in Docker Swarm mode. As described in \SubSec{subsec:04_docker_swarm}, Docker Swarm mode will maintain a healthy state of all running containers. Therefore, the environment will enable self-healing according to the requirements of Autonomic Computing described in \Sec{sec:02_ac}.


\subsection{Apache Spark Cluster}
\label{subsec:05_arch_spark}


\subsubsection{Master and Worker}
The Apache Spark cluster will consist of a Spark master node and a dynamic number of Spark worker nodes.
% Master
The Spark master node is responsible to distribute the application workload across available Spark worker nodes.
% Worker
A Spark worker node will execute the workload given by the master node. Each Spark worker is homogeneous. 
% Homegeneous spark worker
Homegenneosity is important to scale the number of worker nodes. To enable homogeneous nodes, each SPark worker node is a Docker container running the same Docker image. In addition, each worker is given the same computing resources. With homogeneous Spark worker nodes, each worker will respond as all other nodes.
% GPU acceleration
To enable GPU acceleration, each WORKER/MASTER will have the RAPIDS plugin installed.
% Standalone mode
The cluster will be deployed in standalone mode. To be able to run Python applications


\subsubsection{Spark Submit}
Because the Apache Spark cluster ill be executed in standalone mode, a node inside the cluster is required to run Spark applications. When a Spark application will be executed, a Spark Submit container will be deployed in the cluster. When the appliccation has finished, the container will be automatically removed.
Each app will be executed by a unique Spark Submit node.
The Spark Submit node will be deployed via the CI pipeline (SECTION XY). The prupose of this 


\subsection{Autonomic Manager}
\label{subsec:05_arch_am}

\begin{figure}[h]
\centering
\includegraphics[scale=0.85]{images/05_conceptual_design/autonomic_manager/autonomic_manager_overview}
\caption{Autonomic manager component design - Source: Authors own model.}
\label{fig:am-design-component}
\end{figure}

% MAPE architecture
The design of the autonomic manager needs to fullfill all requirements given by the MAPE architecture (described in \SubSec{subsec:02_ac_manager}). It will be responsible to monitor the performance of Apache Spark worker nodes in the environment, analyze the metrics and plan and execute scaling actions in accordance to fullfill the performance goals.
% Figure
As illustrated in \Fig{fig:am-design-component}, the autonomic manager consists of a cAdvisor, Prometheus and Auto-Scaler node. Together, all three nodes build a complete autonomic manager in accordance to the MAPE architecture. Each node will run as a Docker container.
% Explain figure
For monitoring, cAdvisor collects metrics from all available Docker container in the computing environment. Prometheus pulls metrics from cAdvsior. The Auto-Scaler will fetch metrics from Prometheus and send scaling instructions to the Docker engine.


\subsubsection{Monitoring System}
% Why - dynamic env, requirements, etc
The design introduced in \SubSec{subsec:05_arch_overall} is a dynamic changing computing environment. To monitor this dynamic environment, a monitoring-system is needed that fullfills the requirements described in \Sec{sec:02_monitoring}.


% Gathering metrics
cAdvisors will be used as an agent to collect performance metrics from all running Docker containers in the environment.
% Store metrics
Prometheus pulls the collected performance metrics from cAdvisor and stores the data as time-series data in its database.
% Query
In addition, Prometheus provides a powerful multi-dimensional query language to aggregate and analyze the stored data.


\subsubsection{Workflow}

\paragraph{} The workflow of the autonomic manager is implemented as a loop.

% Workflow
\begin{figure}[h]
\centering
\includegraphics[scale=0.50]{images/05_conceptual_design/autonomic_manager/autonomic_manager_workflow}
\caption{UML activity model of the autonomic manager process - Source: Authors own model.}
\label{fig:am-workflow}
\end{figure}
% Explain figure
\paragraph{} \Fig{fig:am-workflow} illustrates all steps of each component of the autonomic manager process according to the MAPE architecture.


% Monitor
The workflow starts with collecting metrics in the monitor phase. cAdvisor is responsible to collect metrics from the Docker engine. After that, Prometheus stores the collected metrics as time-series data in its database.
% Analyze
Next, in the analyze phase, the Auto-Scaler needs to determine if a scaling action is needed. If scaling the Apache Spark worker nodes is not needed, the process has finished and will repeat from the beginning in the next period.
% Plan
If the performance is over- or under-utilized, a scaling action is needed. Then, the Auto-Scaler needs to determine how many Apache Spark worker nodes are needed to reach the performance goal.
% Execute
Lastly, the Auto-Scaler is responsible to send the scaling instructions to the Docker engine.


% ===========================================
% ===========================================
\section{Auto-Scaler}
\label{sec:05_auto-scaler}
% What is the Auto-Scaler
The Auto-Scaler is a component of the the Autonomic Manager and is responsible for the analyze, plan and execution phase.
% Complete autonomic manager
Together with cAdvisor and Prometheus, the Auto-Scaler builds a complete Autonomic Manager according to \SubSec{subsec:02_ac_manager}.
% control loop
In addition, the Auto-Scaler implements the control-loop which is responsible to make adjustments in the environment according to the performance.
% Communication with Prometheus
To adjust the number of Apache Spark worker in the environment, the Auto-Scaler needs to perform queries on the performance metrics stored in Prometheus.
% Analyzing
After analyzing the performance metrics, the Auto-Scaler will send instructions to adjust the number of Apache Spark worker to the Docker engine. The Docker engine tells the Swarm manager to scale the replicas of the Apache Spark service in accordance to the Auto-Scaler instructions.

% HIER NOCH BILD VON SCALER WIE MIT DER MAPE ARCH


\subsection{Configuration}
\label{subsec:05_auto-scaler_configuration}
The Auto-Scaler needs specific configuration properties to be able to collect the correct metrics from Prometheus and deploy new Apache Spark worker container in the environment. The following are properties that have to be defined to ensure that the Auto-Scaler is able to collect meaningful metrics and scale Apache Spark worker as expected.

\subsubsection{General properties}

\begin{itemize}
\item \textbf{Interval seconds:} The number of seconds when the loop has to repeat needs to be defined.

\item \textbf{Cooldown period:} The duration in seconds,  the Auto-Scaler has to wait after a scaling action was performed.

\item \textbf{Recurrence factor:} To prevent to many scaling actions,  the autonomic manager should only execute a scaling action,  if the utilization thresholds is violated \textit{n} times.

\item \textbf{Prometheus URL:} The Auto-Scaler will fetch the configured metrics from the Prometheus REST API.
\end{itemize}

\subsubsection{Metrics}

To support to analyze multiple metrics, the user should be able to create a dynamic list if metrics. Each metric needs to have a variety of properties configured.

\begin{itemize}
\item \textbf{Target utilization:} The relative target utilization of a metrics needs to be defined to calculate the number of Spark worker to add or to remove to reach the defined goal.

\item \textbf{Utilization thresholds:} To determine if a scaling action is needed, the scaling heat algorithm needs the minimum and maximum utilization defined by an administrator.

\item \textbf{Query:} A PromQL query needs to be defined to collect the metric for all Spark Worker.
\end{itemize}

\subsubsection{Apache Spark worker properties}

\begin{itemize}
\item \textbf{Worker image:} To guarantee that each Spark worker is homogeneous, all worker container should be created with the same image.

\item \textbf{Worker network:} To establish communication between all Spark worker and the Spark master, all new Spark worker container should be in the same network.

\item \textbf{Worker thresholds:} The minimum and maximum number of concurrent Spark worker should be defined. To avoid the cold start effect, the minimum amount of worker should be 1. 

\item \textbf{Apache Spark master URI:} To distribute the workload across all Spark Worker, all Spark Worker need to communicate with the Spark master.
\end{itemize}


\subsection{Analyze}
% Short intro to analyze phase
In order to determine if a scaling-action is necessary, the Auto-Scaler has to process the collected metrics. 
% How to analyze
During each period, the Auto-scaler queries the Prometheus time-series database with the configured queries to get all needed metrics. 
% Determine if scaling is needed
After the metrics are received, the Auto-Scaler determines if a scaling action is needed using the Scaling Heat algorithm (introduced in Section AB). If scaling is not necessary, the Auto-Scaler continues to collect metrics from Prometheus.


\subsection{Plan}
% Short intro
If a scaling-action is necessary, the Auto-Scaler is responsible to plan how to scale the number of Spark worker to satisfy the defined utilization goals.
% What is a scaling plan
A scaling plan consists of instructions to add or remove Spark worker which will be send to the Docker engine.
% Calculating Spark worker
To calculate the number of Spark worker, needed to accomplish the defined target utilization, the Auto-Scaler uses the \textit{Kubernetes Horizontal Pod Auto-Scaling} algorithm. In addition, the Auto-Scaler needs to check if the estimated number of Spark worker fall bellow the minimum threshold or exceed the maximum threshold of concurrent Spark worker.


\subsection{Execute}
% Short intro
After a scaling plan has been created, the Auto-Scaler needs to send the instructions to the Docker engine.
% Cooldown
After scaling the environment, it needs time for changes to take effect. Therefore a cooldown period will be activated after each scaling action.
% Whats happens
During the cooldown period, no scaling actions will be forwarded to the Docker engine.

\chapter{Implementation}
\label{sec:implementation}
%
\todo{Describe Chapter}

\section{General}
% BLa bla alles in Docker usw.
T create an elastic computing environment, all needed components will be run as Docker container in a Docker network.


\section{Computing environment}
% The environment will be implemented with Docker/docker-compose
To create an elastic environment Docker will be used. With docker Spark worker can be easy scaled. With docker-compose an environment can be created.
% Images used
% network (Communication)
\subsection{Monitoring}
cAdvisor will be uses to monitor metrics from the Docker engine. Prometheus collects the metrics defined in Section XY from cAdvisor and stores them in its time-series database.
% The queries
The query used for getting CPU metrics:

\begin{lstlisting}[language=SQL, caption=Python example]
SUM(SELECT BLA FROM XYZ)
\end{lstlisting}

The query used for getting GPU metrics:


\subsection{Apache Spark Images}
The are 4 different Spark images:

1. Spark-Base
2. Spark-Master
3. Spark-Worker
4. Spark-Submit

Spark-Base is the foundation image for all other Spark images. Spark-Master is the image for the master node. Spark-Worker is the image for all SPark worker nodes. With Spark-Submit, a application will be submitted, the container runs as long as the application runs and exits after the execution automatically. Spark-Submit is used, because Cluster runs in Standalaone mode. Python cann submit cannot run in cluster mode QUELLE SPark.

\subsection{GPU Acceleration}

NVIDIA RAPIDS will be used as the plugin for GPU acceleration.



\section{Auto-Scaler}
The Auto-Scaler consists of a Python module which implements the logic for control-loop and a own Docker image.
% Is implemented in python
The Auto-Scaler is implemented in the Python programming language.
% Is its own Docker image
The Auto-Scaler module will be executed in an individual Docker image.
% UML diagram

\subsection{Configuration}
All configuration properties defined in the concept of the Auto-Scaler in SECTION A will be defined in a YAML file.
% As a cmd-line argument
To set the configuration properties, the configuration file needs to be set as a command-line argument.
% File example

\subsection{Control-Loop}
Control-Loop of the Auto-Scaler implements the Analyze, Plan, and Execute phase of the MAPE architecture.

\subsubsection{Analyze}
% Scaling Heat
Scaling Heat algorithm will be used to estimate if a scaling action is necessary.

\subsubsection{Plan}
% KHPA
The KHPA algorithm will be used to calculate how many worker are needed to reach the target utilization. Since we have two different metrics, the highest number of running worker will be used KUBERNETES QUELLE.

\subsubsection{Execute}
% Python Docker SDK
Via the Python Docker SDK new container will be spawned in the network.
% Check if worker is busy
If worker need to be removed, it is necessary to check if the worker are running any applications at the moment. Spark provides a REST API to check this.
% Cooldown
After a scaling action has been performed, a cooldown period will be applied so all nodes can relax.



\section{Automated Deployment}
Hier die CI/CD von gitlab.

\chapter{Evaluation}
\label{chap:07_evaluation}

In this chapter, the implementation of \Chap{chap:06_implementation} is evaluated in various experiments. Further, the experimental results are summarised and discussed.

\section{Experimental Environment}
%
The experiments are conducted on the DGX introduced in \Sec{sec:06_impl-env}.
% 2 GPUs
As for the implementation, 2 \hyperlink{abbr:gpu}{GPUs} have been available for the experiments.


% ===========================================
% ===========================================
\section{Experiments}
% Intro
The performance of the computing environment is measured on two widely used Machine Learning algorithms.
% Explain
In both experiments, a Machine Learning model is trained on the computing environment using Apache Spark. Furthermore, each benchmark is evaluated using three different configurations:
\begin{enumerate}
\item Using a static number of \hyperlink{abbr:cpu}{CPU}-only Apache Spark workers (\Sec{sec:07_static})
\item Using a static number of \hyperlink{abbr:gpu}{GPU}-accelerated Apache Spark workers (\Sec{sec:07_gpu})
\item Dynamically scaling the replicas of \hyperlink{abbr:cpu}{CPU}-only Apache Spark workers using the \textit{Auto-Scaler} (\Sec{sec:07_auto-scaler})
\end{enumerate}
% 10 times
For each configuration, the experiment is conducted ten times. The spark-job execution time mean value over all ten iterations per experiment is used as the result.
% DIfference between GPU and CPU
The performance difference between \hyperlink{abbr:gpu}{GPU}-accelerated worker nodes, and \hyperlink{abbr:cpu}{CPU}-only worker nodes are explored using a \hyperlink{abbr:cpu}{CPU}-only and a \hyperlink{abbr:gpu}{GPU}-only version of the benchmark implementation.

\paragraph{}
% Why no auto gpu
The \textit{Auto-Scaler} is not evaluated using \hyperlink{abbr:gpu}{GPU}-accelerated worker nodes. As mentioned, for this experiment, only 2 \hyperlink{abbr:gpu}{GPUs} are available. The RAPIDS plugin allocates a \hyperlink{abbr:gpu}{GPU} exclusively for each executor. When the \textit{Auto-Scaler} creates a new worker node, the worker allocates a random available GPU.
% Live system
However, the host machine is a live system. It is not possible to allocate a random \hyperlink{abbr:gpu}{GPU} that may be in use by another application. Furthermore, the \textit{Auto-Scaler} does not support allocating a specific \hyperlink{abbr:gpu}{GPU} for each newly created worker. An approach to overcome this problem is introduced in \Sec{subsec:08_outlook_gpus}.

\paragraph{}
% The algorithm
The two benchmarks to evaluate the performance are:
\begin{itemize}
\item \hyperlink{abbr:xgboost}{XGBoost} classification model using the \textit{Fannie Mae's Single-Family Historical Loan Performance Dataset}\footnote{Downloaded from: \url{https://docs.rapids.ai/datasets/mortgage-data} (Accessed: 2021-02-06)}\cite{Fannie2021Mortgage}

\item \hyperlink{abbr:xgboost}{XGBoost} regression model using a Taxi fare dataset\footnote{The Taxi dataset is available at: \url{https://github.com/NVIDIA/spark-xgboost-examples/tree/spark-3} (Accessed: 2021-02-06)}
\end{itemize}
% Source
The source code and the dataset used in these experiments are available on Github on the \textit{spark-xgboost-examples}\footnote{spark-xgboost-examples - \url{https://github.com/NVIDIA/spark-xgboost-examples/tree/spark-3} (Accessed: 2021-02-06)} repository from NVIDIA.
% The two diff impl conf
The repository provides a \textit{mortgage} and a \textit{taxi} application. Both applications provide a \hyperlink{abbr:cpu}{CPU} and a \hyperlink{abbr:gpu}{GPU} configuration.

% supervised
Classification and regression algorithms are supervised Machine Learning algorithms.
% Classification
The classification algorithm identifies the category of a label.
% Regression
A regression algorithm predicts a continuous numeric value.
% whats supervised
The goal of supervised Machine Learning algorithms is to train a model by finding patterns in labeled data. The trained model is then used to predict labels on newly observed data based on the learned labels \cite{Mcdonald2020SparkRapids}.


% ===========================================
% ===========================================
\section{Static CPU-Only Workers}
\label{sec:07_static}
% Intro
The first experiment is conducted using a fixed number of \hyperlink{abbr:cpu}{CPU}-only Apache Spark workers to evaluate both benchmarks' performance.
% Appendix
The performance data of each iteration for both benchmarks is available at \Sec{sec:appendix_eval}.


% results
\begin{figure}[h]
\centering
\includegraphics[scale=0.9]{images/07_evaluation/overall_cpu}
\caption{Mean execution time using CPU-only worker nodes}
\label{fig:07_static_results}
\end{figure}
% Explain fig
\Fig{fig:07_static_results} illustrates the mean execution time of both benchmarks for different numbers of worker nodes.
% classification
The best performance for the classification benchmark was achieved with 247 seconds using 30 workers with an improvement of 7.04x compared to 1 worker.
% regression
The regression benchmark performed best with 102 seconds using 10 workers with an improvement of 3.21x in comparison to 1 worker. It achieved a mean execution time of 102 seconds with 30 workers as well. However, using 10 workers is considered the best result for the regression benchmark because fewer workers are used.

\paragraph{}
% The improvement curve
In addition, \Fig{fig:07_static_results} shows that from a specific number of workers onwards, no significant improvement is made.
% Classification
For the classification benchmark, this point is reached at 15 workers.
% Regression
A similar result can be seen for the regression benchmark. From 10 workers onward, the execution time does not improve anymore.
% and?
This result shows that over-provisioning the Apache Spark cluster by simply adding more workers does not improve the performance.


% ===========================================
% ===========================================
\section{GPU-Accelerated Workers}
\label{sec:07_gpu}
% Intro
The impact of \hyperlink{abbr:gpu}{GPU}-accelerated Apache Spark worker nodes for both benchmarks have been evaluated using 2 configurations: Using one and two \hyperlink{abbr:gpu}{GPU}-accelerated worker nodes respectively.
%
The performance data of this experiment is available at \Sec{sec:appendix_eval}.


% vs figure
\begin{figure}[h]
\centering
\includegraphics[scale=0.9]{images/07_evaluation/overall_cpu_vs_gpu}
\caption{Mean execution time of GPU-accelerated worker nodes vs. Mean execution time of CPU-only worker nodes}
\label{fig:07_gpu_results}
\end{figure}
% Explain results
The results of the experiment are illustrated in \Fig{fig:07_gpu_results}.
% vs
The mean execution times using \hyperlink{abbr:gpu}{GPU}-accelerated worker nodes are plotted against the mean execution times using \hyperlink{abbr:cpu}{CPU}-only worker nodes.
% overall
Overall, \hyperlink{abbr:gpu}{GPU}-accelerated worker nodes significantly outperformed the equivalent number of \hyperlink{abbr:cpu}{CPU}-only worker nodes.
% classification
For the classification benchmark, the best execution time was achieved using 2 \hyperlink{abbr:gpu}{GPU}-accelerated worker nodes with an improvement of 2.42x compared to 2 \hyperlink{abbr:cpu}{CPU}-only worker nodes. Using 1 \hyperlink{abbr:gpu}{GPU}-accelerated worker node achieved an improvement of 2.91x compared to 1 \hyperlink{abbr:cpu}{CPU}-only worker node.
% regression
The regression benchmark performed best with 2 \hyperlink{abbr:gpu}{GPU}-accelerated worker nodes improving the execution by 3.02x in comparison to 2 \hyperlink{abbr:cpu}{CPU}-only worker nodes. Using 1 \hyperlink{abbr:gpu}{GPU}-accelerated worker node achieved an improvement of 4.55x in comparison to 1 \hyperlink{abbr:cpu}{CPU}-only worker node.

\paragraph{}
% The table
\Tab{table:07_gpu_overall_results} shows the overall results by comparing \hyperlink{abbr:cpu}{CPU}-only and \hyperlink{abbr:gpu}{GPU}-accelerated results for each benchmark.
% classification
For the classification benchmark, 5 \hyperlink{abbr:cpu}{CPU}-only worker nodes outperformed 1 \hyperlink{abbr:gpu}{GPU}-accelerated worker node, and 10 \hyperlink{abbr:cpu}{CPU}-only worker nodes performed better than 2 \hyperlink{abbr:gpu}{GPU}-accelerated worker nodes.
% regression
However, for the regression benchmark, \hyperlink{abbr:cpu}{CPU}-only worker nodes could not outperform \hyperlink{abbr:gpu}{GPU}-accelerated worker node performance.
% Results as table
\begin{table}[ht]
\centering
\begin{tabular}{@{}l|ll|ll@{}}
\toprule
                  & \multicolumn{2}{c|}{Classification}                & \multicolumn{2}{c}{Regression}                    \\
Number of workers & \multicolumn{1}{c}{CPU} & \multicolumn{1}{c|}{GPU} & \multicolumn{1}{c}{CPU} & \multicolumn{1}{c}{GPU} \\ \midrule
1  & 1740s         & 597s          & 328s          & 72s \\
2  & 828s          & \textbf{342s} & 139s          & \textbf{46s} \\
5  & 535s          & -             & 129s          & -     \\
10 & 295s          & -             & \textbf{102s} & -     \\
15 & 263s          & -             & 109s          & -     \\
20 & 281s          & -             & 120s          & -     \\
25 & 257s          & -             & 129s          & -     \\
30 & \textbf{247s} & -             & 102s          & -     \\
35 & 258s          & -             & 109s          & -     \\ \bottomrule
\end{tabular}
\caption{Mean execution time of all CPU-only and GPU-accelerated experiments for both benchmarks}
\label{table:07_gpu_overall_results}
\end{table}
\newpage


% ===========================================
% ===========================================
\section{Auto-Scaler}
\label{sec:07_auto-scaler}
% Intro
To test the impact of the \textit{Auto-Scaler} while training Machine Learning applications, both benchmarks have been tested with different \textit{Auto-Scaler} configurations.
% Performance metric
The performance data of both benchmarks is available at \Sec{sec:appendix_eval}.


\subsection{Auto-Scaler Benchmark Configurations}
% Wheres it from
The configuration parameters are chosen according to the results of the static worker experiment (\Sec{sec:07_static}).
% Table
The \textit{Auto-Scaler} configuration parameters for each benchmark are shown in \Tab{table:07_auto-scaler_config_parameter}.
% Iteration Figure
\begin{figure}[h]
\centering
\includegraphics[scale=0.9]{images/07_evaluation/overall_auto-scaler_iterations}
\caption{CPU utilization during all iterations of both benchmarks with 2 CPU-only workers}
\label{fig:07_auto-scaler_iterations_results}
\end{figure}
% Max worker
To prevent cluster over-provisioning, the maximum number of worker nodes for the classification benchmark is set to 15, and 10 for the regression benchmark respectively.
% recurrence
\Fig{fig:07_auto-scaler_iterations_results} shows all 10 iterations of the \hyperlink{abbr:cpu}{CPU}-only experiment for both benchmarks. It shows that, while the Machine Learning model is trained, only 2 performance spikes occurred at the beginning and at the end during a short period. To scale while these performance spikes occur, the recurrence factor is set to 1. Additionally, the maximum \hyperlink{abbr:cpu}{CPU} utilization for both benchmarks is set according to the maximum \hyperlink{abbr:cpu}{CPU} utilization during the performance spikes.
% The configurations
\begin{table}[ht]
\centering
\begin{tabular}{@{}l|ll@{}}
\toprule
Parameter               & Classification & Regression \\ \midrule
Interval                & 5 seconds      & 5 seconds  \\
Recurrence factor       & 1              & 1          \\
Cooldown period         & 60 seconds     & 60 seconds \\
Target CPU utilization  & 10\%           & 5\%        \\
Minimum CPU utilization & 5\%            & 2\%       \\
Maximum CPU utilization & 25\%           & 10\%       \\
Minimum worker nodes    & 2              & 2         \\
Maximum worker nodes    & 15             & 10         \\ \bottomrule
\end{tabular}
\caption{\textit{Auto-Scaler} configuration parameters for both benchmarks}
\label{table:07_auto-scaler_config_parameter}
\end{table}


\subsection{Execution Time Evaluation}
% Figure
\Fig{fig:07_auto-scaler_results} illustrates the results of the \textit{Auto-Scaler} experiment. The mean execution time of 2 \hyperlink{abbr:cpu}{CPU}-only worker nodes is plotted against the mean execution with the \textit{Auto-Scaler} enabled.
% Explain
The results show that no improvement was achieved using the \textit{Auto-Scaler} implementation.
% Results
For the classification benchmark, the execution time increased by 228 seconds when enabling the \textit{Auto-Scaler} in comparison to 2 \hyperlink{abbr:cpu}{CPU}-only worker nodes. Furthermore, the execution time is 38 seconds higher with \textit{Auto-Scaler} enabled than 2 \hyperlink{abbr:cpu}{CPU}-only worker nodes.
% Figure
\begin{figure}[h]
\centering
\includegraphics[scale=0.9]{images/07_evaluation/overall_auto-scaler}
\caption{\textit{Auto-Scaler} experiment mean execution time of both benchmarks}
\label{fig:07_auto-scaler_results}
\end{figure}

\newpage
\paragraph{}
% Table
\begin{table}[ht]
\centering
\begin{tabular}{@{}l|ll|ll@{}}
\toprule
                               & \multicolumn{2}{c|}{Classification}                   & \multicolumn{2}{c}{Regression}                       \\
\multicolumn{1}{c|}{Iteration} & \multicolumn{1}{c}{Time} & \multicolumn{1}{c|}{Nodes} & \multicolumn{1}{c}{Time} & \multicolumn{1}{c}{Nodes} \\ \midrule
1  & \textbf{840s} & 15 & 168s          & 5 \\
2  & 960s          & 7  & \textbf{150s} & 2 \\
3  & 1020s         & 15 & 186s          & 6 \\
4  & 1140s         & 15 & 180s          & 6 \\
5  & 1080s         & 15 & 180s          & 6 \\
6  & 1140s         & 15 & 180s          & 6 \\
7  & 1140s         & 15 & 180s          & 6 \\
8  & 1080s         & 15 & 180s          & 6 \\
9  & 1140s         & 7  & 180s          & 6 \\
10 & 1020s         & 15 & 186s          & 6 \\ \bottomrule
\end{tabular}
\caption{Results of each iteration during the \textit{Auto-Scaler} experiment for both benchmarks}
\label{table:07_auto-scaler_iterations}
\end{table}
% Explain
\Tab{table:07_auto-scaler_iterations} summarizes the results for all 10 iterations of each benchmark.
% classification
For the classification benchmark, the best result was achieved in the first iteration with an execution time of 960 seconds and a maximum number of 15 worker nodes during the execution.
% regression
The regression benchmark performed best in the second iteration, with an execution time of 150 seconds and a maximum number of 2 worker nodes during this iteration.
% n maximum reached
The results show that, during the classification benchmark, the \textit{Auto-Scaler} scaled upon the maximum number of allowed worker nodes in 8 out of 10 iterations. During the regression benchmark experiment, the maximum number of worker nodes was never reached. Furthermore, in the second iteration, the \textit{Auto-Scaler} had not scaled the number of worker nodes when the best result was achieved.


\section{Experimental Results Summary}
% Intro
Overall, three different experiments have been conducted to evaluate the implementation of this thesis.
% First experiment
The results of the first experiment show that over-provisioning the Apache Spark cluster does not improve the performance significantly and can negatively impact the performance. 
%  Second
The second experiment shows that \hyperlink{abbr:gpu}{GPU}-accelerated worker nodes significantly outperformed \hyperlink{abbr:cpu}{CPU}-only worker nodes. However, \hyperlink{abbr:cpu}{CPU}-only workers achieved a better performance in the classification benchmark than \hyperlink{abbr:gpu}{GPU}-accelerated workers with 5x the amount.
% Third
For the last experiment, the results show that enabling the \textit{Auto-Scaler} during the training of Machine Learning models increases the execution time compared with a fixed number of worker nodes.

\chapter{Conclusion and Outlook}
\label{chap:08_outlook}

In this chapter, an outlook is given for further research based on the results of \Chap{chap:07_evaluation}.


% ===========================================
% ===========================================
\section{Conclusion}
% Intro
In this thesis, a scalable computing environment for training machine learning models using Apache Spark was implemented. In addition, the training process is automated using an automated deployment pipeline.
% The research questions
In \Chap{chap:01_introduction}, three research questions for this thesis are defined.


% RQ1
The first research question is about implementing a solution to automatically scale the number of workers in an Apache Spark cluster.
% The solution
For this research question, a computing environment according to the autonomic computing architecture is implemented.
% Autonomic computing
It consists of an autonomic manager and an Apache Spark cluster.
% Docker
For simplicity, all components of the computing environment were deployed with Docker.
% manager
The autonomic manager is implemented according to the MAPE architecture and responsible for monitoring and automatically scale the Apache Spark workers.
% monitoring system
For the implementation of the autonomic manager, Prometheus, cadvisor, and DCGM-Exporter are used to create a monitoring system.
% Auto-Scaler
In addition, an \textit{Auto-Scaler} is implemented using Python 3 that automatically fetches performance metrics from the monitoring system and scales the Apache Spark worker nodes.


% RQ2
The second research question is about how to accelerate the Apache Spark cluster's computational power using GPUs.
% RAPIDS
This research question is implemented by extending the existing Apache Spark cluster with the NVIDIA RAPIDS plugin.


% RQ3
The final research question is about how to automate the training of a machine learning model.
%
To accommodate this research question, a CI pipeline is implemented using GitLab CI/CD. The CI pipeline is performed when a change is committed to the machine learning applications' source code.
% How
It consists of a \textit{train-stage} that automatically deploys a \textit{spark-submit} Docker container to the Apache Spark cluster.
% spark-submit
The \textit{spark-submit} container executes the spark-submit executable to perform the application on the Apache Spark cluster.
% Gitlab runner
To deploy a \textit{spark-submit} container from the \textit{train-stage} inside the Apache Spark cluster, a GitLab Runner is created, which runs on the same host.


% Results
The implementation is evaluated with a variety of experiments using two widely used machine learning algorithms.
% GPU
The GPU-accelerated Apache Spark cluster significantly outperformed a CPU-only cluster.
% AUto-Scaler
However, the \textit{Auto-Scaler} implementation is not able to improve the performance of the Apache Spark cluster. Furthermore, it increases the overall execution time.


% ===========================================
% ===========================================
\section{Outlook}
% Intro
The \textit{Auto-Scaler} implementation was not able to improve the performance of training machine learning models on the Apache Spark cluster.
% A possible reason
A possible reason for the execution time increase is that Apache Spark is not able to efficiently distribute the workload of an application while workers are added to the cluster during the execution of an application.
% Differet approaches
Therefore, different \textit{Auto-Scaler} implementation approaches for further research are introduced.


\subsection{Proactive Auto-Scaler Approach}
% Currently
The current implementation of the \textit{Auto-Scaler} follows a reactive approach. It uses static threshold-based rules (e.g., maximum and minimum CPU utilization) to define when a scaling action is necessary.
% Not sufficient
The \textit{Auto-Scaler} can be optimized using a proactive approach. A proactive \textit{Auto-Scaler} tries to predict the computing environment's future state (e.g., reinforcement learning) using historical data.
% Extansion of the MAPE architecture
The MAPE architecture can be extended with a knowledge source (MAPE-K), which provides data about the computing environment \cite{Sinreich2006AnAB}. The autonomic manager can use the knowledge to try to predict when a scaling action is necessary.


\subsection{Vertical-Scaling Approach}
% evaluations
The evaluation of the static CPU-only workers' experiment showed that from a specific number of workers, no significant performance improvement is made.
% vertical
Therefore, instead of scaling the replicas of worker nodes, a vertical-scaling approach, where the \textit{Auto-Scaler} scales the available computing resources of a worker node, should be explored.
% results
Providing more powerful resources increases the computational power of the Apache Spark cluster and should lead to lower execution times for training machine learning models.


\subsection{Scaling GPU-accelerated Workers}
\label{subsec:08_outlook_gpus}
% Currently no GPUs
Due to the lack of implementation, the \textit{Auto-Scaler} can not keep track of available GPUs in the cluster. Therefore, it was not possible to evaluate the \textit{Auto-Scalers} impact on GPU-accelerated worker nodes.
% GPU approach
Extending the \textit{Auto-Scaler} to keep track of GPUs which are free or already allocated by a worker should be explored.
%
Auto-scaling GPU-accelerated worker nodes can improve the spark-job execution time.


\subsection{Save Temporary Data on an External Service}
% The problem
With the \textit{Auto-Scaler's} current implementation, it is not possible to remove worker nodes while an application is actively performing because worker nodes save important temporary data.
% A solution
The worker nodes can be extended to save temporary data on an external storage service, which enables the \textit{Auto-Scaler} to remove nodes while an application is performed.
% 
The impact of removing nodes during the execution of an application should be explored. In the evaluation of this thesis, only the impact of adding worker nodes is explored.

%\fi

% Anhang (Bibliographie darf im deutschen nicht in den Anhang!)
\bibliography{literatur}
%\nocite{*}
\clearpage

% Anhang
\appendix
% 'Anhang' ins Inhaltsverzeichnis
\phantomsection
\addcontentsline{toc}{part}{Appendix}
\chapter{Attached CD}
\label{chap:appendix_cd}
The source code of this implementation is available on the CD attached to this thesis. It includes the following directories:
\begin{itemize}
\item auto-scaler: Source code of the \textit{Auto-Scaler} implementation
\item computing\_framework: Source code of the computing environment
\item spark-images: Source code of the Apache Spark images
\item ci-pipeline: Source code of the CI pipeline
\end{itemize}

\chapter{Implementation}

\section{Computing Environment}
% docker compose file
\begin{lstlisting}[label=lst:appendix_env_compose, caption=Computing environment docker-compose file]
version: "3.7"
 
networks:
    computing-net:
        name: computing_net
        attachable: true
 
services:
    spark-master:
        image: spark-master:3.0.1-hadoop2.7
        networks:
            - computing-net
        ports:
            - 4040:4040
            - 7077:7077
        volumes:
            - ./conf/spark-master/metrics.properties:/usr/bin/spark/conf/metrics.properties
            - ./conf/spark-master/spark-defaults.conf:/usr/bin/spark/conf/spark-defaults.conf
 
    prometheus:
        image: prom/prometheus
        networks:
            - computing-net
        volumes:
            - ./conf/prometheus/prometheus.yml:/etc/prometheus/prometheus.yml
            - ./conf/prometheus/recording_rules.yml:/etc/prometheus/recording_rules.yml
        command:
            - "--config.file=/etc/prometheus/prometheus.yml"
        ports:
            - "9090:9090"
        depends_on:
            - cadvisor
 
    cadvisor:
        image: google/cadvisor
        networks:
            - computing-net
        ports:
            - "8080:8080"
        volumes:
            - "/:/rootfs:ro"
            - "/var/run:/var/run:ro"
            - "/sys:/sys:ro"
            - "/var/lib/docker/:/var/lib/docker:ro"
            - "/dev/disk/:/dev/disk:ro"
        command:
            - "--docker_only=true"
            - "--logtostderr=true"
        depends_on:
            - spark-master
            - spark-worker
 
    auto-scaler:
        image: auto-scaler:latest
        networks:
            - computing-net
        volumes:
            - "/var/run/docker.sock:/var/run/docker.sock:ro"
            - ./conf/auto-scaler/config.yml:/etc/autoscaler/config.yml
        command:
            - '--config=/etc/autoscaler/config.yml'
        depends_on:
            - prometheus
\end{lstlisting}

\chapter{Apache Spark Cluster Implementation}
% Base image dockerfile
\begin{lstlisting}[label=lst:appendix_spark_base_dockerfile, caption=Apache Spark base image Dockerfile]
FROM nvidia/cuda:11.0-devel-ubuntu16.04
 
LABEL maintainer="marcel.pascal.stolin@ipa.fraunhofer.de"
 
ARG SPARK_VERSION
ARG HADOOP_VERSION
 
# Install all important packages
RUN apt-get update -qy && \
    apt-get install -y openjdk-8-jre-headless procps python3 python3-pip curl
 
# Install Apache Spark
RUN mkdir /usr/bin/spark/ && \
    curl https://ftp-stud.hs-esslingen.de/pub/Mirrors/ftp.apache.org/dist/spark/spark-${SPARK_VERSION}/spark-${SPARK_VERSION}-bin-hadoop${HADOOP_VERSION}.tgz -o spark.tgz && \
    tar -xf spark.tgz && \
    mv spark-${SPARK_VERSION}-bin-hadoop${HADOOP_VERSION}/* /usr/bin/spark/ && \
    rm -rf spark.tgz && \
    rm -rf spark-${SPARK_VERSION}-bin-hadoop${HADOOP_VERSION}/
 
# Add GPU discovery script
RUN mkdir /opt/sparkRapidsPlugin/
COPY getGpusResources.sh /opt/sparkRapidsPlugin/getGpusResources.sh
ENV SPARK_RAPIDS_DIR=/opt/sparkRapidsPlugin
 
# Install cuDF and RAPIDS
RUN curl -o ${SPARK_RAPIDS_DIR}/cudf-0.15-cuda11.jar https://repo1.maven.org/maven2/ai/rapids/cudf/0.15/cudf-0.15-cuda11.jar
RUN curl -o ${SPARK_RAPIDS_DIR}/rapids-4-spark_2.12-0.2.0.jar https://repo1.maven.org/maven2/com/nvidia/rapids-4-spark_2.12/0.2.0/rapids-4-spark_2.12-0.2.0.jar
ENV SPARK_CUDF_JAR=${SPARK_RAPIDS_DIR}/cudf-0.15-cuda11.jar
ENV SPARK_RAPIDS_PLUGIN_JAR=${SPARK_RAPIDS_DIR}/rapids-4-spark_2.12-0.2.0.jar
 
# Set all environment variables
ENV JAVA_HOME=/usr/lib/jvm/java-8-openjdk-amd64
ENV SPARK_HOME /usr/bin/spark
ENV SPARK_NO_DAEMONIZE true
ENV PYSPARK_DRIVER_PYTHON python3
ENV PYSPARK_PYTHON python3
ENV PATH /usr/bin/spark/bin:/usr/bin/spark/sbin:$PATH
 
WORKDIR ${SPARK_HOME}
\end{lstlisting}

% Master image dockerfile
\begin{lstlisting}[label=lst:appendix_spark_master_dockerfile, caption=Apache Spark master image Dockerfile]
ARG SPARK_VERSION
ARG HADOOP_VERSION
 
FROM spark-base:$SPARK_VERSION-hadoop$HADOOP_VERSION
 
LABEL maintainer="marcel.pascal.stolin@ipa.fraunhofer.de"
 
# Set ports
ENV SPARK_MASTER_PORT 7077
ENV SPARK_MASTER_WEBUI_PORT 4040
 
EXPOSE 4040 7077
 
# Start master-node in standalone mode
ENTRYPOINT [ "sbin/start-master.sh" ]
\end{lstlisting}

% Worker image Dockerfile
\begin{lstlisting}[label=lst:appendix_spark_worker_dockerfile, caption=Apache Spark worker image Dockerfile]
ARG SPARK_VERSION
ARG HADOOP_VERSION
 
FROM spark-base:$SPARK_VERSION-hadoop$HADOOP_VERSION
 
LABEL maintainer="marcel.pascal.stolin@ipa.fraunhofer.de"
 
# Add spark-env
COPY spark-env.sh ${SPARK_HOME}/conf/spark-env.sh
 
# Set port
ENV SPARK_WORKER_WEBUI_PORT 4041
 
EXPOSE 4041
 
# Start worker-node
ENTRYPOINT ./sbin/start-slave.sh ${SPARK_MASTER_URI}
\end{lstlisting}

% GPU discovery script
\begin{lstlisting}[label=lst:appendix_spark_gpu-discovery, caption=GPU discovery script, language=bash]
#!/usr/bin/env bash

#
# Licensed to the Apache Software Foundation (ASF) under one or more
# contributor license agreements.  See the NOTICE file distributed with
# this work for additional information regarding copyright ownership.
# The ASF licenses this file to You under the Apache License, Version 2.0
# (the "License"); you may not use this file except in compliance with
# the License.  You may obtain a copy of the License at
#
#    http://www.apache.org/licenses/LICENSE-2.0
#
# Unless required by applicable law or agreed to in writing, software
# distributed under the License is distributed on an "AS IS" BASIS,
# WITHOUT WARRANTIES OR CONDITIONS OF ANY KIND, either express or implied.
# See the License for the specific language governing permissions and
# limitations under the License.
#

# This script is a basic example script to get resource information about NVIDIA GPUs.
# It assumes the drivers are properly installed and the nvidia-smi command is available.
# It is not guaranteed to work on all setups so please test and customize as needed
# for your environment. It can be passed into SPARK via the config
# spark.{driver/executor}.resource.gpu.discoveryScript to allow the driver or executor to discover
# the GPUs it was allocated. It assumes you are running within an isolated container where the
# GPUs are allocated exclusively to that driver or executor.
# It outputs a JSON formatted string that is expected by the
# spark.{driver/executor}.resource.gpu.discoveryScript config.
#
# Example output: {"name": "gpu", "addresses":["0","1","2","3","4","5","6","7"]}

ADDRS=`nvidia-smi --query-gpu=index --format=csv,noheader | sed -e ':a' -e 'N' -e'$!ba' -e 's/\n/","/g'`
echo {\"name\": \"gpu\", \"addresses\":[\"$ADDRS\"]}
\end{lstlisting}

% Custom submit script
\begin{lstlisting}[label=lst:appendix_spark-submit_script, caption=Custom submit script, language=sh]
#!/bin/bash
 
DRIVER_MEMORY=${DRIVER_MEMORY-4g}
EXECUTOR_MEMORY=${EXECUTOR_MEMORY-8g}
 
if [ "$CPU_ONLY" == "true" ]
    then
    echo "Submit Spark app with CPU only"
 
    $SPARK_HOME/bin/spark-submit \
        --master $SPARK_MASTER_URI \
        --driver-memory $DRIVER_MEMORY \
        --executor-memory $EXECUTOR_MEMORY \
        --conf spark.driver.extraClassPath=${SPARK_CUDF_JAR}:${JAR_RAPIDS}:${LIBS_PATH}/xgboost4j_3.0-1.3.0-0.1.0.jar:${LIBS_PATH}/xgboost4j-spark_3.0-1.3.0-0.1.0.jar \
        --conf spark.executor.extraClassPath=${SPARK_CUDF_JAR}:${JAR_RAPIDS}:${LIBS_PATH}/xgboost4j_3.0-1.3.0-0.1.0.jar:${LIBS_PATH}/xgboost4j-spark_3.0-1.3.0-0.1.0.jar \
        --jars ${SPARK_CUDF_JAR},${LIBS_PATH}/xgboost4j-spark_3.0-1.3.0-0.1.0.jar,${LIBS_PATH}/xgboost4j_3.0-1.3.0-0.1.0.jar,${JAR_RAPIDS} \
        --py-files ${LIBS_PATH}/xgboost4j-spark_3.0-1.3.0-0.1.0.jar,/tank/data/users/chh-ms/spark-xgboost-examples/examples/apps/python/samples.zip \
        $@
else
    echo "Submit Spark app with GPU acceleration"
 
    RAPIDS_GPU_ALLOC_FRACTION=${RAPIDS_GPU_ALLOC_FRACTION-1}
    RAPIDS_INCOMPATIBLE_OPS=${RAPIDS_INCOMPATIBLE_OPS-"false"}
    RAPIDS_DEBUG=${RAPIDS_DEBUG-"NONE"}
    RAPIDS_GPU_POOL=${RAPIDS_GPU_POOL-"ARENA"}
 
    $SPARK_HOME/bin/spark-submit \
        --master $SPARK_MASTER_URI \
        --driver-memory $DRIVER_MEMORY \
        --executor-memory $EXECUTOR_MEMORY \
        --conf spark.plugins=com.nvidia.spark.SQLPlugin \
        --conf spark.rapids.memory.gpu.pool=$RAPIDS_GPU_POOL \
        --conf spark.rapids.memory.gpu.allocFraction=$RAPIDS_GPU_ALLOC_FRACTION \
        --conf spark.rapids.sql.incompatibleOps.enabled=$RAPIDS_INCOMPATIBLE_OPS \
        --conf spark.rapids.memory.gpu.debug=$RAPIDS_DEBUG \
        --conf spark.task.resource.gpu.amount=1 \
        --conf spark.executor.resource.gpu.amount=1 \
        --conf spark.driver.extraClassPath=${SPARK_CUDF_JAR}:${JAR_RAPIDS}:${LIBS_PATH}/xgboost4j_3.0-1.3.0-0.1.0.jar:${LIBS_PATH}/xgboost4j-spark_3.0-1.3.0-0.1.0.jar \
        --conf spark.executor.extraClassPath=${SPARK_CUDF_JAR}:${JAR_RAPIDS}:${LIBS_PATH}/xgboost4j_3.0-1.3.0-0.1.0.jar:${LIBS_PATH}/xgboost4j-spark_3.0-1.3.0-0.1.0.jar \
        --jars ${SPARK_CUDF_JAR},${LIBS_PATH}/xgboost4j-spark_3.0-1.3.0-0.1.0.jar,${LIBS_PATH}/xgboost4j_3.0-1.3.0-0.1.0.jar,${JAR_RAPIDS} \
        --py-files ${LIBS_PATH}/xgboost4j-spark_3.0-1.3.0-0.1.0.jar,/tank/data/users/chh-ms/spark-xgboost-examples/examples/apps/python/samples.zip \
        $@  
fi
\end{lstlisting}

% Custom submit script
\begin{lstlisting}[label=lst:appendix_spark-build-script, caption=Build script for all Apache Spark images, language=sh]
#!/bin/bash
 
if [ $# -ge 2 ]
  then
    SPARK_VERSION=$1
    HADOOP_VERSION=$2
 
    PWD=$(pwd)
 
    IMAGE_TAG_PATH="local"
 
    echo
    echo "# Image tag-path: ${IMAGE_TAG_PATH}"
    echo
 
    echo "# STEP(1/4) - Building spark-base"
    docker build \
      -t $IMAGE_TAG_PATH/spark-base:$SPARK_VERSION-hadoop$HADOOP_VERSION \
      --build-arg SPARK_VERSION=$SPARK_VERSION \
      --build-arg HADOOP_VERSION=$HADOOP_VERSION \
      $PWD/spark-base/
 
    echo "# STEP(2/4) - Building spark-master"
    docker build \
      -t $IMAGE_TAG_PATH/spark-master:$SPARK_VERSION-hadoop$HADOOP_VERSION \
      --build-arg IMAGE_TAG_PATH=$IMAGE_TAG_PATH \
      --build-arg SPARK_VERSION=$SPARK_VERSION \
      --build-arg HADOOP_VERSION=$HADOOP_VERSION \
      $PWD/spark-master/
 
    echo "# STEP(3/4) - Building spark-worker"
    docker build \
      -t $IMAGE_TAG_PATH/spark-worker:$SPARK_VERSION-hadoop$HADOOP_VERSION \
      --build-arg IMAGE_TAG_PATH=$IMAGE_TAG_PATH \
      --build-arg SPARK_VERSION=$SPARK_VERSION \
      --build-arg HADOOP_VERSION=$HADOOP_VERSION \
      $PWD/spark-worker/
 
    echo "# STEP(4/4) - Building spark-submit"
    docker build \
      -t $IMAGE_TAG_PATH/spark-submit:$SPARK_VERSION-hadoop$HADOOP_VERSION \
      --build-arg IMAGE_TAG_PATH=$IMAGE_TAG_PATH \
      --build-arg SPARK_VERSION=$SPARK_VERSION \
      --build-arg HADOOP_VERSION=$HADOOP_VERSION \
      $PWD/spark-submit/
  else
    echo "No arguments supplied\n"
    echo "Use the script as follows:"
    echo "build-images.sh <SPARK_VERSION> <HADOOP_VERSION>"
    exit 1
fi
\end{lstlisting}


\section{GitLab CI/CD Pipeline Implementation}

% Custom submit script
\begin{lstlisting}[label=lst:appendix_ci-pipeline_job, caption=texttt{.gitlab-ci.yml} configuration file]
stages:
    - test
    - train
 
unit-tests:
    stage: test
    image: python:3.8-slim
    tags:
        - spark-submit
    script:
        - pip3 install -e .
        - cd tests
        - pytest
 
train-model:
    stage: train
    tags:
        - spark-submit
    variables:
        TMP_MOUNT_POINT: /tmp/_mnt
        NETWORK: "computing_net"
        CPU_ONLY: "true"
        DRIVER_MEMORY: "4g"
        EXECUTOR_MEMORY: "8g"
        SPARK_MASTER_URI: "spark://spark-master:7077"
    script:
        - pwd
        - mkdir -p "$TMP_MOUNT_POINT"
        - mv -v * $TMP_MOUNT_POINT
        - 'export MOUNT_POINT="$(dirname $CI_PROJECT_DIR)/shared"'
       - rm -rf $MOUNT_POINT
        - mkdir -p "$MOUNT_POINT"
        - cp -R $TMP_MOUNT_POINT/* $MOUNT_POINT/
        - rm -rf $TMP_MOUNT_POINT
        - chmod -R 777 $MOUNT_POINT
        - sh $MOUNT_POINT/submit.sh /mnt/train.py --max-depth=50
    only:
        - master
\end{lstlisting}


\section{Evaluation Data}
\label{sec:appendix_eval}

\begin{figure}[h]
\centering
\includegraphics[scale=0.5]{images/appendix/evaluation_data/classification_benchmark/classification_1_worker_cpu_performance}
\caption{Classification benchmark - Performance data of all iterations with 1 worker}
\label{fig:appendix_eval_classification_static1}
\end{figure}

\begin{figure}[h]
\centering
\includegraphics[scale=0.5]{images/appendix/evaluation_data/classification_benchmark/classification_2_worker_cpu_performance}
\caption{Classification benchmark - Performance data of all iterations with 2 worker}
\label{fig:appendix_eval_classification_static2}
\end{figure}

\begin{figure}[h]
\centering
\includegraphics[scale=0.5]{images/appendix/evaluation_data/classification_benchmark/classification_5_worker_cpu_performance}
\caption{Classification benchmark - Performance data of all iterations with 5 worker}
\label{fig:appendix_eval_classification_static5}
\end{figure}

\begin{figure}[h]
\centering
\includegraphics[scale=0.5]{images/appendix/evaluation_data/classification_benchmark/classification_10_worker_cpu_performance}
\caption{Classification benchmark - Performance data of all iterations with 10 worker}
\label{fig:appendix_eval_classification_static10}
\end{figure}

\begin{figure}[h]
\centering
\includegraphics[scale=0.5]{images/appendix/evaluation_data/classification_benchmark/classification_15_worker_cpu_performance}
\caption{Classification benchmark - Performance data of all iterations with 15 worker}
\label{fig:appendix_eval_classification_static15}
\end{figure}

\begin{figure}[h]
\centering
\includegraphics[scale=0.5]{images/appendix/evaluation_data/classification_benchmark/classification_20_worker_cpu_performance}
\caption{Classification benchmark - Performance data of all iterations with 20 worker}
\label{fig:appendix_eval_classification_static20}
\end{figure}

\begin{figure}[h]
\centering
\includegraphics[scale=0.5]{images/appendix/evaluation_data/classification_benchmark/classification_25_worker_cpu_performance}
\caption{Classification benchmark - Performance data of all iterations with 25 worker}
\label{fig:appendix_eval_classification_static25}
\end{figure}

\begin{figure}[h]
\centering
\includegraphics[scale=0.5]{images/appendix/evaluation_data/classification_benchmark/classification_30_worker_cpu_performance}
\caption{Classification benchmark - Performance data of all iterations with 30 worker}
\label{fig:appendix_eval_classification_static30}
\end{figure}

\begin{figure}[h]
\centering
\includegraphics[scale=0.5]{images/appendix/evaluation_data/classification_benchmark/classification_35_worker_cpu_performance}
\caption{Classification benchmark - Performance data of all iterations with 35 worker}
\label{fig:appendix_eval_classification_static35}
\end{figure}

\begin{figure}[h]
\centering
\includegraphics[scale=0.5]{images/appendix/evaluation_data/classification_benchmark/classification_gpu1_performance}
\caption{Classification benchmark - Performance data of all iterations with 1 GPU-accelerated worker}
\label{fig:appendix_eval_classification_gpu1}
\end{figure}

\begin{figure}[h]
\centering
\includegraphics[scale=0.5]{images/appendix/evaluation_data/classification_benchmark/classification_gpu2_performance}
\caption{Classification benchmark - Performance data of all iterations with 2 GPU-accelerated workers}
\label{fig:appendix_eval_classification_gpu2}
\end{figure}

\begin{figure}[h]
\centering
\includegraphics[scale=0.5]{images/appendix/evaluation_data/classification_benchmark/classification_auto-scaler_performance}
\caption{Classification benchmark - Performance data of all iterations with the \textit{Auto-Scaler} enabled}
\label{fig:appendix_eval_classification_auto-scaler}
\end{figure}

\begin{figure}[h]
\centering
\includegraphics[scale=0.5]{images/appendix/evaluation_data/regression_benchmark/regression_1_worker_cpu_performance}
\caption{Regression benchmark - Performance data of all iterations with 1 worker}
\label{fig:appendix_eval_regression_static1}
\end{figure}

\begin{figure}[h]
\centering
\includegraphics[scale=0.5]{images/appendix/evaluation_data/regression_benchmark/regression_2_worker_cpu_performance}
\caption{Regression benchmark - Performance data of all iterations with 2 worker}
\label{fig:appendix_eval_regression_static2}
\end{figure}

\begin{figure}[h]
\centering
\includegraphics[scale=0.5]{images/appendix/evaluation_data/regression_benchmark/regression_5_worker_cpu_performance}
\caption{Regression benchmark - Performance data of all iterations with 5 worker}
\label{fig:appendix_eval_regression_static5}
\end{figure}

\begin{figure}[h]
\centering
\includegraphics[scale=0.5]{images/appendix/evaluation_data/regression_benchmark/regression_10_worker_cpu_performance}
\caption{Regression benchmark - Performance data of all iterations with 10 worker}
\label{fig:appendix_eval_regression_static10}
\end{figure}

\begin{figure}[h]
\centering
\includegraphics[scale=0.5]{images/appendix/evaluation_data/regression_benchmark/regression_15_worker_cpu_performance}
\caption{Regression benchmark - Performance data of all iterations with 15 worker}
\label{fig:appendix_eval_regression_static15}
\end{figure}

\begin{figure}[h]
\centering
\includegraphics[scale=0.5]{images/appendix/evaluation_data/regression_benchmark/regression_20_worker_cpu_performance}
\caption{Regression benchmark - Performance data of all iterations with 20 worker}
\label{fig:appendix_eval_regression_static20}
\end{figure}

\begin{figure}[h]
\centering
\includegraphics[scale=0.5]{images/appendix/evaluation_data/regression_benchmark/regression_25_worker_cpu_performance}
\caption{Regression benchmark - Performance data of all iterations with 25 worker}
\label{fig:appendix_eval_regression_static25}
\end{figure}

\begin{figure}[h]
\centering
\includegraphics[scale=0.5]{images/appendix/evaluation_data/regression_benchmark/regression_30_worker_cpu_performance}
\caption{Regression benchmark - Performance data of all iterations with 30 worker}
\label{fig:appendix_eval_regression_static30}
\end{figure}

\begin{figure}[h]
\centering
\includegraphics[scale=0.5]{images/appendix/evaluation_data/regression_benchmark/regression_35_worker_cpu_performance}
\caption{Regression benchmark - Performance data of all iterations with 35 worker}
\label{fig:appendix_eval_regression_static35}
\end{figure}

\begin{figure}[h]
\centering
\includegraphics[scale=0.5]{images/appendix/evaluation_data/regression_benchmark/regression_gpu1_performance}
\caption{Regression benchmark - Performance data of all iterations with 1 GPU-accelerated worker}
\label{fig:appendix_eval_regression_gpu1}
\end{figure}

\begin{figure}[h]
\centering
\includegraphics[scale=0.5]{images/appendix/evaluation_data/regression_benchmark/regression_gpu2_performance}
\caption{Regression benchmark - Performance data of all iterations with 2 GPU-accelerated workers}
\label{fig:appendix_eval_regression_gpu2}
\end{figure}

\begin{figure}[h]
\centering
\includegraphics[scale=0.5]{images/appendix/evaluation_data/regression_benchmark/regression_auto-scaler_performance}
\caption{Regression benchmark - Performance data of all iterations with the \textit{Auto-Scaler} enabled}
\label{fig:appendix_eval_regression_auto-scaler}
\end{figure}


\IfDefined{printindex}{\printindex}
\IfDefined{printnomenclature}{\printnomenclature}



%% Dokument ENDE %%%%%%%%%%%%%%%%%%%%%%%%%%%%%%%%%%%%%%%%%%%%%%%%%%%%%%%%%%
\end{document}

